\chapter{Fragilidad bancaria y riesgos del crecimiento a la baja: evidencia empírica del spread soberano}
\label{chap:paper2}

\section*{Resumen del capítulo}
\addcontentsline{toc}{section}{Resumen del capítulo}

Este capítulo evalúa empíricamente si la fragilidad sistémica del sector bancario y el riesgo de la cola izquierda del crecimiento del producto determinan de forma \emph{complementaria} —no aditiva— el riesgo de crédito soberano en economías emergentes. La fragilidad bancaria se mide mediante la métrica \textit{Joint Loss} ($JLoss$) \citep{Chari2024b}, un indicador semiparamétrico a nivel país de la pérdida conjunta esperada del sistema bancario, construido a partir de probabilidades de incumplimiento de tipo \citet{Merton1974b} agregadas vía aproximación de punto de silla \citep{MartinThompsonBrowne2001}; los balances y precios accionarios de 113 bancos que alimentan la métrica se extraen de forma homogénea de la terminal Bloomberg. El riesgo a la baja del crecimiento se captura mediante \textit{Growth-at-Risk} ($GaR$), estimado por regresión cuantílica de panel \citep{Koenker2004} siguiendo el marco de \textit{vulnerable growth} de \citet{ABG2019b} y la implementación de CEMLA \citep{OssandonBusch2022}. El spread soberano se mide \emph{exclusivamente} por el CDS soberano a cinco años (Bloomberg). Sobre \textbf{un único panel} de catorce economías emergentes con $JLoss$, $GaR$ y CDS disponibles simultáneamente (2004Q1--2026Q1; $N=838$, y $N=738$ y trece países bajo la especificación con el vector completo de controles domésticos), el coeficiente de interacción $JLoss_{i,t}\times GaR_{i,t}$ es negativo, del signo predicho por el multiplicador de \textit{doom loop} ($\hat\theta=-0{,}35$, $t=-1{,}91$): significativo al 10\% bajo errores de Driscoll--Kraay ($p=0{,}056$) y bajo \textit{wild cluster bootstrap} ($p=0{,}035$), y al 1\% bajo \textit{clustering} por país ($p=0{,}001$). Un modelo de umbral de panel \citep{Hansen1999,Hansen2000} corrobora la no linealidad: el efecto de la fragilidad sobre el spread es $+8{,}1$ puntos básicos por unidad de $JLoss$ en el régimen de cola severa frente a $+2{,}3$ en el benigno, y el efecto marginal de la fragilidad crece monótonamente a medida que el riesgo de cola se agrava. El hallazgo se reporta con dos fronteras explícitas: es un fenómeno \textbf{posterior a la crisis financiera global} —negativo y significativo en todas las ventanas móviles de cinco años que empiezan en 2012, nulo o de signo contrario en 2004--2011, sin quiebre discreto en 2020— y su significancia descansa de forma desproporcionada en el puñado de países con variación sustantiva de $JLoss$ (al excluir China deja de ser significativa). Una estrategia de variables instrumentales \textit{shift-share} con un instrumento \textbf{fuerte} (primera etapa $F=21{,}6$; exposición pre-muestral estimada antes de 2012) respalda con solidez la dirección banco$\to$soberano del canal de nivel ($\hat\beta_{JLoss}=+17{,}5$ pb, $p=0{,}001$); un segundo instrumento (choque cambiario del dólar) tiene primera etapa igual de fuerte pero falla el test de sobre-identificación al combinarse con el primero, una advertencia que se reporta con franqueza. Este capítulo constituye la arista empírica de la tesis; el Capítulo~\ref{chap:paper1}, de naturaleza teórica, deriva un modelo de organización industrial bancaria del cual esta complementariedad es una predicción de primer orden, y pone a prueba —con resultado negativo sobre este panel— su amplificación por la concentración del sistema bancario.

\section{Introducción}

Durante las crisis financieras de las últimas dos décadas los spreads soberanos de las economías emergentes se han ampliado de forma más pronunciada que los de las economías avanzadas. Una parte de esa ampliación es atribuible a la fragilidad del sector bancario doméstico —que eleva el pasivo público contingente y deteriora los balances soberanos por la vía del rescate— y, de manera decisiva, a su coincidencia con un debilitamiento agudo de las perspectivas de crecimiento. La conjetura que organiza este capítulo es que esos dos canales no operan de forma aditiva: la fragilidad bancaria y el riesgo a la baja del crecimiento se \textbf{complementan} en su transmisión al riesgo soberano, de modo que la coincidencia de ambos amplifica el spread por encima de la suma de sus efectos individuales.

La literatura relevante ha tratado esos dos canales por separado. La del nexo banca--soberano ha documentado con detalle cómo la fragilidad del sistema bancario se transmite al riesgo soberano, pero trata el crecimiento como un fundamento de fondo o como un factor de confusión a controlar. La literatura de \textit{growth-at-risk} ha caracterizado con precisión cómo las condiciones financieras deterioran la cola izquierda de la distribución del crecimiento, pero detiene su análisis en el producto sin proyectarlo sobre la prima soberana. El trabajo más próximo, \citet{Chari2024b}, introduce una métrica estructural de fragilidad bancaria sistémica en la determinación del spread soberano de economías emergentes y explora su interacción con factores financieros \textbf{globales} —no con una medida doméstica y endógena del riesgo de cola del crecimiento. La pregunta de investigación es, por tanto: ¿amplifica la coincidencia de alta fragilidad bancaria y elevado riesgo a la baja del crecimiento el riesgo soberano por encima de la suma de sus efectos individuales?

Para responderla, este capítulo integra dos métricas del estado del arte —la fragilidad bancaria sistémica medida por la métrica \textit{Joint Loss} ($JLoss$) de \citet{Chari2024b} y el riesgo de cola del crecimiento medido por \textit{Growth-at-Risk} ($GaR$) en la tradición de \citet{ABG2019b}— en un único marco de determinación del spread soberano, y caracteriza la \textbf{no linealidad} de su interacción mediante un término multiplicativo y un modelo de umbral. La evidencia se construye sobre un \textit{pipeline} reproducible de extracción de balances bancarios y de estimación de $GaR$ en pseudo--tiempo real, sobre un único panel de catorce economías emergentes cuyos insumos de mercado provienen íntegramente de la terminal Bloomberg. El aporte no es la mera adición de un regresor, sino la articulación —teórica y empírica— de dos canales que la literatura ha mantenido separados, reportada con una delimitación explícita de dónde el mecanismo está establecido y dónde sigue siendo una hipótesis abierta.

Este capítulo es la arista empírica de una tesis compuesta por dos investigaciones sobre el mismo fenómeno. El Capítulo~\ref{chap:paper1}, de naturaleza teórica, desarrolla un modelo de organización industrial bancaria en el que la estructura de mercado —la intensidad de competencia entre bancos— es el parámetro profundo que gobierna la cadena que va desde la fragilidad sistémica hasta el riesgo soberano, y del cual la complementariedad $JLoss\times GaR$ que se estima aquí es una predicción de primer orden. Ese capítulo pone a prueba, además, la predicción que distingue al modelo de la organización industrial bancaria estándar \citep{MMR2010b}: que la complementariedad se amplifica en los sistemas bancarios más concentrados. Por esa razón, la triple interacción $JLoss\times GaR\times HHI$ se reporta en el Capítulo~\ref{chap:paper1} y no aquí.

El resto del capítulo se organiza así. La Sección~2 desarrolla el marco teórico y deriva las hipótesis contrastables; la Sección~3 sitúa el trabajo en la literatura; la Sección~4 detalla la construcción de $JLoss$ y $GaR$ y la estrategia econométrica; la Sección~5 describe el panel y su cobertura; la Sección~6 reporta los resultados, la robustez y la identificación causal; y la Sección~7 discute las implicancias y concluye. El Anexo~\ref{chap:anexoB} tabula la procedencia de todas las series.

\section{Marco teórico}

El argumento central de este trabajo se sostiene sobre tres pilares teóricos que la literatura ha desarrollado de forma mayoritariamente independiente. El primero es el bucle de retroalimentación soberano-bancario (\textit{doom loop}), que articula la transmisión bidireccional de fragilidad entre los balances del sistema bancario y los del soberano. El segundo es la interpretación de la fragilidad bancaria sistémica como un pasivo contingente del soberano, que provee el fundamento microeconómico de la métrica $JLoss$. El tercero es la determinación del riesgo a la baja del crecimiento por las condiciones financieras —el programa de investigación de \textit{vulnerable growth}—, que sustenta la métrica $GaR$.

La contribución conceptual de este capítulo no reside en ninguno de estos pilares por separado, sino en la tesis de que operan de manera \textbf{complementaria} y no aditiva: la fragilidad bancaria y la vulnerabilidad del crecimiento se amplifican mutuamente en su transmisión al riesgo soberano. Esta sección desarrolla cada pilar y formaliza la predicción de amplificación supra-aditiva que orienta la estrategia empírica de las Secciones 4 a 6.

\subsection{El nexo soberano-bancario: del bucle de retroalimentación al abrazo mortal}

La piedra angular teórica es el mecanismo de retroalimentación entre la solvencia bancaria y la solvencia soberana formalizado por \citet{FarhiTirole2018b}. En su modelo de tres fechas, el soberano emite deuda cuyo precio refleja la capacidad fiscal esperada, mientras los bancos domésticos gestionan su liquidez manteniendo bonos soberanos —domésticos y extranjeros— como reserva de valor.

Ante un choque adverso en la fecha intermedia, la solvencia del soberano se deteriora por dos vías: un efecto \textbf{directo}, en tanto el choque contrae la actividad económica y, con ella, los ingresos fiscales; y un efecto \textbf{indirecto}, que opera a través de los balances bancarios. La caída inicial del precio de la deuda pública erosiona el patrimonio neto de los bancos y deprime su inversión; si el soberano valora suficientemente esa inversión, rescata al sistema bancario. El punto decisivo es que el rescate se financia emitiendo deuda adicional, lo que reduce todavía más su precio: por cada dólar de rescate, el acervo de deuda pública debe aumentar en más de un dólar.

Formalmente, este mecanismo admite una representación de punto fijo para el precio de la deuda soberana doméstica en la fecha intermedia, $p_1(s) = 1 - F\!\left(B_1(s)\mid s\right)$, donde $B_1(s)$ es el acervo de deuda estado-contingente y $F(\cdot)$ la distribución de la capacidad fiscal. La sensibilidad de $p_1(s)$ al estado $s$ adopta la forma de un \textbf{multiplicador de amplificación}: una caída del precio incrementa el rescate requerido y el volumen de bonos a colocar, lo que vuelve a deprimir el precio, \textit{ad infinitum}. La magnitud de este multiplicador es creciente en el grado de \textit{home bias}, esto es, en la proporción de deuda soberana en manos de bancos domésticos.

La implicación que este trabajo explota es que la amplificación es \textbf{intrínsecamente no lineal}: el deterioro soberano no es la suma de un componente bancario y uno fiscal independientes, sino su producto a través de un multiplicador que se intensifica con la fragilidad del balance bancario y con la severidad del estado adverso.

\subsection{La fragilidad bancaria como pasivo contingente del soberano: el canal \texorpdfstring{$JLoss$}{JLoss}}

La métrica de fragilidad bancaria empleada en este trabajo, $JLoss$, se interpreta como el costo directo esperado de rescatar al conjunto del sistema bancario ante un evento sistémico, y por ello entra naturalmente en la valoración del bono soberano \citep{Chari2024b}. Esta interpretación operacionaliza el ``\textit{bailout put}'' de \citet{FarhiTirole2018b} y la dilución de \citet{AcharyaDrechslerSchnabl2014b}: si el mercado anticipa que un soberano deberá asumir el costo de recapitalizar su banca en un escenario de tensión, ese pasivo contingente debe estar incorporado \textit{ex ante} en la prima soberana.

\subsection{El riesgo a la baja del crecimiento: el canal \texorpdfstring{$GaR$}{GaR}}

El tercer pilar proviene del programa de \textit{vulnerable growth} de \citet{ABG2019b}, cuya tesis es que las condiciones financieras predicen de manera \textbf{asimétrica} la distribución condicional del crecimiento futuro del producto: la dependencia respecto de las condiciones financieras es marcadamente más fuerte en los cuantiles inferiores —la cola izquierda— que en los superiores, que permanecen comparativamente estables en el tiempo. Sus hallazgos asociados —la correlación negativa entre la media y la volatilidad condicionales del crecimiento, y la ``paradoja de la volatilidad'' según la cual períodos de baja volatilidad anteceden a resultados de crecimiento negativos \citep{BrunnermeierSannikov2014,AdrianBoyarchenko2012}— implican que el objeto relevante para la evaluación del riesgo no es el crecimiento esperado, sino la severidad de su escenario adverso plausible.

La traducción de este canal al riesgo soberano opera a través de la dinámica de sostenibilidad de la deuda. La ley de movimiento del ratio de deuda sobre producto, $\dot{d} = (r - g)\,d - pb$ —donde $d$ es la deuda/PIB, $r$ la tasa de interés real, $g$ la tasa de crecimiento y $pb$ el balance primario—, muestra que un sesgo a la baja en la distribución de crecimiento deteriora la trayectoria esperada de sostenibilidad y, sobre todo, la capacidad fiscal precisamente en los estados de cola donde el incumplimiento se torna probable. Dado que el costo del default posee un perfil de pago de tipo opción —discreto y concentrado en el peor escenario—, lo determinante para la prima soberana no es el primer momento de la distribución de crecimiento, sino su cola izquierda. Es esta no linealidad la que conecta el canal de crecimiento con el canal soberano y la que hace insuficiente controlar por el crecimiento medio: se requiere una medida del riesgo de cola, esto es, $GaR$.

Operativamente, $GaR$ se define como el cuantil de cola izquierda de la distribución condicional del crecimiento (con la convención de que valores más negativos indican mayor riesgo a la baja), estimado mediante regresión cuantílica \citep{KoenkerBassett1978,Koenker2004} sobre un índice de condiciones financieras. El concepto, nacido como extensión macroeconómica del \textit{value-at-risk} \citep{PrasadEtAl2019}, se implementa en este trabajo siguiendo la plataforma de código abierto de CEMLA \citep{OssandonBusch2022}, que organiza los factores de riesgo del crecimiento en tres bloques —condiciones macrofinancieras locales, un índice de condiciones financieras y condiciones internacionales ortogonalizadas— y caracteriza la densidad condicional completa del crecimiento en lugar de únicamente su media. El puente conceptual entre condiciones financieras laxas, acumulación de vulnerabilidad y materialización del riesgo de cola lo ofrece \citet{Stein2012}, al concebir la política monetaria como instrumento de regulación de la estabilidad financiera.

\subsection{Síntesis: complementariedad y amplificación supra-aditiva}

Los tres pilares convergen en una predicción común que constituye la hipótesis teórica central del trabajo: \textbf{el efecto de la fragilidad bancaria sobre el riesgo soberano no es independiente del estado del crecimiento, y viceversa}. La coincidencia de alta fragilidad bancaria y elevado riesgo a la baja del crecimiento amplifica el spread soberano por encima de la suma de los efectos individuales de cada dimensión.

La justificación de esta complementariedad se desprende de la articulación de los tres pilares:

\begin{itemize}
    \item \textbf{Del mecanismo de \citet{FarhiTirole2018b}:} el multiplicador de amplificación es creciente tanto en la fragilidad del balance bancario como en la severidad del estado adverso $s$. Identificando $GaR$ con la cola izquierda del crecimiento —es decir, con los estados adversos $s$—, el efecto cruzado del multiplicador implica que la pérdida de precio del bono soberano ante un incremento de la fragilidad es mayor en estados de bajo crecimiento. La amplificación es, por construcción, un fenómeno de interacción.
    \item \textbf{Del canal de dilución de \citet{AcharyaDrechslerSchnabl2014b}:} la sensibilidad del riesgo soberano a las condiciones bancarias escala con el grado de deterioro pre-existente, y el factor macroeconómico común que vincula banca y soberano es el crecimiento. La interacción fragilidad $\times$ estado macro es, por tanto, una predicción directa del modelo.
    \item \textbf{De la complementariedad de \citet{GennaioliMartinRossi2014}:} el costo del incumplimiento soberano crece con la exposición bancaria; un escenario de bajo crecimiento que eleva la probabilidad de default, combinado con un sistema bancario frágil, magnifica el costo total de la cesación de pagos y, con él, la prima soberana.
\end{itemize}

Estos argumentos motivan el objeto empírico central del trabajo: el término de interacción $JLoss_{i,t}\times GaR_{i,t}$ y la predicción de \textbf{complementariedad (super-aditividad)}. Bajo la convención adoptada —$GaR$ medido en niveles como el cuantil de cola izquierda del crecimiento, de modo que valores \textbf{más negativos} señalan mayor riesgo a la baja—, el efecto marginal de la fragilidad bancaria sobre el spread debe intensificarse a medida que el riesgo de cola se agrava. Formalmente, la complementariedad se expresa como

\[ \frac{\partial^{2}\, \mathrm{Spread}_{i,t}}{\partial\, JLoss_{i,t}\;\partial\, GaR_{i,t}} \;<\; 0, \]

esto es, el efecto marginal de la fragilidad bancaria sobre el spread es \textbf{decreciente en $GaR$}: a menor $GaR$ (mayor severidad de la cola), mayor es la sensibilidad del spread a $JLoss$. En la especificación lineal con interacción, $\partial\, \mathrm{Spread}/\partial\, JLoss = \beta_{1} + \theta\,GaR$, por lo que la hipótesis de amplificación se traduce en un coeficiente de interacción \textbf{negativo} ($\theta < 0$): cuando $GaR$ toma valores marcadamente negativos, el término $\theta\,GaR$ es positivo y eleva el efecto de la fragilidad. Una representación de régimen mediante un modelo de umbral \citep{Hansen2000} formaliza la versión no lineal de esta hipótesis: existe un umbral de vulnerabilidad $\lambda$ por debajo del cual (en el régimen de $GaR$ suficientemente adverso) la sensibilidad del spread a la fragilidad bancaria es estructuralmente mayor.

Esta formulación delimita con precisión la contribución del trabajo frente a la literatura previa. \citet{Chari2024b} modelan la interacción de $JLoss$ con factores financieros \textbf{globales} —el VIX, la tasa del Tesoro estadounidense, los spreads de alto rendimiento—, mostrando que los países con banca más frágil son más sensibles a los choques exógenos del ciclo financiero global. El presente trabajo desplaza el foco hacia la interacción de $JLoss$ con una métrica de vulnerabilidad \textbf{doméstica y endógena del crecimiento} ($GaR$), cerrando así el triángulo banca-crecimiento-soberano que la literatura ha analizado únicamente por pares. Las hipótesis contrastables que se derivan del marco son:

\begin{itemize}
    \item \textbf{H1 (canal bancario).} La fragilidad bancaria sistémica se asocia positivamente con el spread soberano: $\beta_{1} > 0$.
    \item \textbf{H2 (canal de crecimiento).} Un mayor riesgo a la baja del crecimiento —un $GaR$ más negativo— se asocia con un mayor spread soberano; en niveles, el signo esperado es $\beta_{2} < 0$.
    \item \textbf{H3 (complementariedad — hipótesis central).} La interacción entre fragilidad bancaria y riesgo a la baja del crecimiento amplifica el spread soberano de manera supra-aditiva, con coeficiente de interacción $\theta < 0$; de forma equivalente, existe un régimen de alta vulnerabilidad ($GaR$ por debajo de un umbral $\lambda$) en el que la sensibilidad del spread a sus determinantes es estructuralmente superior.
\end{itemize}

\section{Estado del arte}

Esta sección sitúa el trabajo en la intersección de cuatro literaturas que han evolucionado de manera predominantemente paralela: la de los determinantes del spread soberano en economías emergentes, la de las métricas de riesgo sistémico y fragilidad bancaria, la del nexo banca-soberano, y la de la frontera entre condiciones financieras y crecimiento del producto (\textit{growth-at-risk}). La revisión persigue un objetivo argumentativo preciso: mostrar que, pese a la madurez de cada literatura por separado, ninguna ha articulado de forma conjunta y no lineal la fragilidad bancaria sistémica y el riesgo a la baja del crecimiento como codeterminantes del riesgo soberano —que es, justamente, la brecha que este trabajo cubre.

\subsection{Determinantes del spread soberano en economías emergentes}

La literatura clásica sobre el riesgo soberano de las economías emergentes ha oscilado entre dos polos explicativos. En un extremo, los \textbf{fundamentos macroeconómicos} domésticos —razón deuda/PIB, crecimiento, términos de intercambio, reservas, calidad institucional— como determinantes del diferencial soberano \citep{Edwards1984,HilscherNosbusch2010}. En el otro, los \textbf{factores globales} o de oferta de capital —el apetito por riesgo, la tasa libre de riesgo estadounidense, la liquidez internacional—, sintetizados en la distinción entre factores \textit{push} y \textit{pull} de la literatura de flujos de capital \citep{CalvoLeidermanReinhart1996}. La evidencia ha tendido a favorecer la preeminencia de los factores globales: \citet{GonzalezRozadaLevyYeyati2008} muestran que una fracción sustancial de la variación de los spreads de los mercados emergentes se explica por el apetito de riesgo global y los spreads de alto rendimiento de las economías avanzadas; \citet{LongstaffEtAl2011} documentan que el riesgo de crédito soberano es mayoritariamente sistemático y se halla más vinculado a los mercados financieros estadounidenses que a los fundamentos locales; y \citet{RemolonaScatignaWu2008} descomponen el spread en una pérdida esperada y una prima de riesgo, atribuyendo la dinámica de esta última a la aversión global al riesgo. La conexión entre el ciclo económico doméstico y el spread fue formalizada por \citet{UribeYue2006}, mientras que la consolidación de un único factor financiero global que gobierna la covariación de los precios de los activos riesgosos a escala mundial —el \textit{Global Financial Cycle}— se debe a \citet{MirandaAgrippinoRey2022}. \citet{CavalloValenzuela2010} y \citet{CsontoIvaschenko2013}, por su parte, descomponen el papel relativo de fundamentos y factores globales en la determinación de los diferenciales. Este trabajo se inserta en esa tradición pero introduce un determinante doméstico hasta ahora ausente del núcleo de la literatura —el riesgo a la baja del crecimiento medido como cuantil de cola— e indaga su interacción con la fragilidad bancaria, en lugar de tratar a ambos como regresores aditivos independientes.

\subsection{Métricas de riesgo sistémico y fragilidad bancaria}

La crisis financiera global de 2008-09 catalizó una vasta literatura orientada a medir el riesgo sistémico del sector financiero. Pueden distinguirse tres familias metodológicas. La primera, basada en datos de mercado y en la teoría de carteras, incluye el \textit{Marginal Expected Shortfall} y el \textit{SRISK} —la insuficiencia de capital condicional del sistema ante un evento sistémico— de \citet{AcharyaEtAl2017b} y \citet{BrownleesEngle2016}, así como el \textit{CoVaR} y su variación $\Delta$CoVaR de \citet{AdrianBrunnermeier2016b}, que cuantifican el riesgo de cola del sistema condicional a la dificultad de una institución individual. La segunda familia, de naturaleza estructural y semiparamétrica, deriva de los modelos de \citet{Merton1974b} y agrega las pérdidas potenciales del sistema mediante técnicas de distribución de pérdidas de cartera: el \textit{Distress Insurance Premium} de \citet{HuangZhouZhu2009} —la prima hipotética de un seguro contra pérdidas sistémicas— y las \textit{Banking Stability Measures} de \citet{SegovianoGoodhart2009} son sus exponentes más cercanos a la métrica empleada en este trabajo. La tercera familia, de corte econométrico, mide la interconexión y la propagación, como en los índices de conectividad de \citet{BillioEtAl2012}; una panorámica de este instrumental se encuentra en \citet{BisiasEtAl2012}.

La métrica $JLoss$ pertenece a la segunda familia y, en rigor, es pariente directa del \textit{Distress Insurance Premium} de \citet{HuangZhouZhu2009}: ambas construyen la distribución agregada de pérdidas del sistema a partir de probabilidades de incumplimiento de tipo Merton y de la dependencia respecto de un factor sistémico. Su especificidad —y su ventaja para el problema soberano— reside en tres rasgos: la agregación a \textbf{nivel de país} (la unidad pertinente para el riesgo soberano y no, como en buena parte de la literatura anglosajona, a nivel de institución o de sistema estadounidense), la \textbf{interpretación como costo directo de rescate} y, decisivamente, la \textbf{disponibilidad para una muestra amplia de economías emergentes}, donde la cobertura de SRISK, CoVaR y medidas afines es escasa o inexistente. La capacidad predictiva comparada de estas métricas sobre la actividad real fue evaluada por \citet{GiglioKellyPruitt2016}, cuyo conjunto de medidas sirve además como batería de controles de robustez en la literatura aplicada y en el presente trabajo.

\subsection{El nexo banca-soberano}

La interdependencia entre la solvencia bancaria y la soberana constituye el corazón teórico del capítulo y ha generado una literatura propia, revitalizada por la crisis de deuda europea. En el plano histórico-empírico, \citet{ReinhartRogoff2009,ReinhartRogoff2011} documentan la recurrencia con que las crisis bancarias anteceden y desencadenan crisis soberanas, y \citet{LaevenValencia2013,LaevenValencia2020} cuantifican los cuantiosos costos fiscales de las crisis bancarias sistémicas en su base de datos de referencia. En el plano teórico, las contribuciones seminales ya discutidas en el marco —\citet{FarhiTirole2018b}, \citet{AcharyaDrechslerSchnabl2014b}, \citet{GennaioliMartinRossi2014}— se complementan con los modelos de riesgo soberano e intermediación de \citet{Bocola2016}, el análisis de fragilidad bancaria en mercados financieramente integrados de \citet{BoltonJeanne2011}, el rol de las garantías gubernamentales en la retroalimentación de \citet{Leonello2018} y la propuesta de activos seguros europeos (ESBies) de \citet{BrunnermeierEtAl2016} como dispositivo para romper el bucle.

En el plano empírico aplicado, \citet{ModySandri2012} muestran que bancos y soberanos quedaron "unidos por la cadera" durante la crisis europea; \citet{DieckmannPlank2012} documentan la transferencia de riesgo del sector privado al público y su modulación por la importancia pre-crisis del sistema financiero doméstico; \citet{KallestrupLandoMurgoci2016} evidencian que los vínculos del sector financiero contienen información sobre el CDS soberano más allá de los fundamentos; y \citet{FratzscherRieth2019} analizan la interacción entre política monetaria, rescates bancarios y el nexo soberano-bancario. La dimensión de la ``\textit{carry trade}'' —bancos que cargan deuda soberana riesgosa de su propio país— fue documentada por \citet{AcharyaSteffen2015}. Esta literatura, sin embargo, es abrumadoramente europea y de frecuencia de crisis; su extensión sistemática a un panel amplio de economías emergentes, con una métrica de fragilidad construida homogéneamente para todos los países, es precisamente la línea que abren \citet{Chari2024b} y que este trabajo profundiza al incorporar el canal de crecimiento. Un modelo estructural que deriva esta misma cadena desde la organización industrial del sector bancario —la estructura de mercado como determinante profundo de la fragilidad de cola y, por esa vía, del spread soberano— se desarrolla en el Capítulo~\ref{chap:paper1}.

\subsection{\textit{Growth-at-Risk} y la frontera finanzas-crecimiento}

La cuarta literatura, de aparición más reciente, vincula el estado de las condiciones financieras con la distribución condicional —y no sólo con la media— del crecimiento futuro del producto. Su contribución fundacional es \citet{ABG2019b}, que establece la asimetría esencial del fenómeno: las condiciones financieras predicen con fuerza la cola izquierda del crecimiento, mientras la cola derecha permanece estable. El marco fue adoptado por los organismos multilaterales como instrumento de vigilancia macroprudencial —el \textit{Global Financial Stability Report} del FMI de 2017 y la aplicación de \citet{PrasadEtAl2019}— y trasladado a las economías latinoamericanas mediante la plataforma de código abierto de CEMLA \citep{OssandonBusch2022}, base metodológica directa de la estimación de $GaR$ de este trabajo. Antecedentes conceptuales del nexo finanzas-cola del crecimiento se hallan en \citet{AllenBaliTang2012}, que indagan si una medida agregada de riesgo de cola del sistema financiero (CATFIN) predice contracciones macroeconómicas, y en \citet{GiglioKellyPruitt2016}. El sustento teórico del canal proviene de la literatura de fijación de precios de activos por intermediarios financieros —\citet{BrunnermeierSannikov2014}, \citet{HeKrishnamurthy2013}—, que racionaliza por qué las fricciones financieras generan no linealidades y riesgo de cola en la actividad real.

Desde el punto de vista econométrico, el método descansa en la regresión cuantílica \citep{KoenkerBassett1978}, sistematizada en \citet{Koenker2005}. Su extensión a datos de panel con efectos fijos —el estimador empleado en este trabajo— plantea desafíos específicos abordados por \citet{Koenker2004}, con penalización de los efectos individuales, y por las alternativas de \citet{Canay2011} y \citet{MachadoSantosSilva2019}; la monotonicidad de la función cuantílica estimada se garantiza mediante el reordenamiento de \citet{ChernozhukovEtAl2010}. La construcción del índice de condiciones financieras que alimenta la estimación se apoya en la tradición de \citet{HatziusEtAl2010} y en indicadores compuestos de estrés sistémico como el CISS de \citet{HolloKremerLoDuca2012}. Es de notar que esta literatura ha mantenido su variable de interés en el crecimiento del producto; su conexión con el \textit{pricing} del riesgo soberano permanece, hasta donde alcanza esta revisión, esencialmente inexplorada.

\subsection{Brecha en la literatura y posicionamiento del trabajo}

La revisión precedente revela una compartimentación persistente. La literatura del nexo banca-soberano (3.3) ha establecido que la fragilidad bancaria eleva el riesgo soberano, pero ha tratado el crecimiento como un fundamento de fondo o como un factor de confusión a controlar —no como un canal de cola con dinámica propia—. La literatura de \textit{growth-at-risk} (3.4) ha caracterizado con precisión cómo las condiciones financieras deterioran la cola izquierda del crecimiento, pero ha detenido su análisis en el producto, sin proyectarlo sobre la prima soberana. Y la literatura de determinantes del spread (3.1) ha privilegiado los factores globales y los fundamentos en niveles, omitiendo las métricas de riesgo de cola del crecimiento. El trabajo más próximo, \citet{Chari2024b}, introduce la fragilidad bancaria sistémica ($JLoss$) en la determinación del spread soberano de economías emergentes y muestra que los países con banca más frágil son más sensibles a los choques \textbf{exógenos} del ciclo financiero global (el VIX, la tasa del Tesoro estadounidense, los spreads de alto rendimiento). La pregunta que este capítulo añade es distinta y complementaria: si esa amplificación también opera a través de una vulnerabilidad \textbf{doméstica y endógena} —el riesgo de cola del propio crecimiento del país—, que no es un choque llegado de fuera sino un estado que la fragilidad bancaria contribuye a generar (canal de \textit{credit crunch}) y que retroalimenta el costo fiscal del rescate. La distinción importa para la política: la primera interacción indica cuándo un sistema bancario frágil expone al soberano al ciclo financiero global; la segunda, cuándo lo expone a su propio ciclo doméstico.

La contribución de este capítulo consiste en cerrar el triángulo banca-crecimiento-soberano. Frente a esa literatura fragmentada, el trabajo (i) emplea conjuntamente una métrica estructural de fragilidad bancaria sistémica ($JLoss$) y una métrica de riesgo de cola del crecimiento ($GaR$) como determinantes del spread soberano; (ii) contrasta la naturaleza \textbf{no lineal} de su interacción —la hipótesis de complementariedad supra-aditiva— mediante un término de interacción y un modelo de umbral \citep{Hansen2000}; y (iii) somete el resultado a una batería de estrategias de identificación causal que declara honestamente los límites de la inferencia con un número reducido de países. En ese sentido, el aporte no es la mera adición de un regresor, sino la articulación teórica y empírica de dos canales que la literatura ha mantenido separados, reportada con el nivel de honestidad estadística que exige extender la conclusión más allá de la muestra donde fue establecida.

\section{Metodología}

La estrategia metodológica se articula en tres bloques. Los dos primeros describen la construcción de las variables explicativas de interés —la fragilidad bancaria sistémica $JLoss_{i,t}$ (Sección 4.1) y el riesgo a la baja del crecimiento $GaR_{i,t}$ (Sección 4.2)—, ambas estimadas a partir de modelos estructurales y semiparamétricos sobre microdatos bancarios y macrofinancieros. El tercer bloque (Sección 4.3) detalla la estrategia de identificación econométrica con la que se contrastan las hipótesis H1–H3. Las dos métricas se construyen de manera independiente y se consolidan en un panel a frecuencia trimestral mediante combinación por país y trimestre.

\subsection{Medición de la fragilidad bancaria: la métrica \texorpdfstring{$JLoss$}{JLoss}}

$JLoss_{i,t}$ es una métrica semiparamétrica, a nivel de país, de la pérdida conjunta esperada del sistema bancario condicional a un evento sistémico, expresada como porcentaje de la exposición total. Su construcción procede en cuatro etapas: (i) la estimación de las probabilidades individuales de incumplimiento de cada banco mediante un modelo estructural de tipo Merton; (ii) la condicionalización de esas probabilidades respecto de un factor sistémico país; (iii) la agregación de la distribución de pérdidas del sistema por el método de punto de silla; y (iv) el cálculo de las contribuciones marginales y la normalización que define la métrica. Se sigue la metodología de \citet{Chari2024b}. Todos los insumos micro —hojas de balance bancarias trimestrales y capitalización bursátil diaria— se obtienen de la terminal Bloomberg mediante un mismo protocolo de extracción para los 113 bancos cotizados de las veinte economías del universo objetivo, lo que elimina la heterogeneidad de definiciones contables entre supervisores nacionales que afectaba a reconstrucciones previas basadas en fuentes regulatorias. La regla de madurez del punto de incumplimiento (deuda emitida = largo plazo; el resto del pasivo = corto plazo) se aplica de manera homogénea a partir de la partida ``\textit{long-term borrowings}'' de Bloomberg.

\subsubsection{Probabilidades individuales de incumplimiento}

Para cada banco se estima la probabilidad de incumplimiento siguiendo la modificación de \citet{Kealhofer2000} del modelo estructural de \citet{Merton1974b}, en su variante KMV. A partir de los pasivos de corto y largo plazo ($L_{ST}$, $L_{LT}$), se define el \textbf{punto de incumplimiento}

\[ D^{*} = L_{ST} + \tfrac{1}{2}\,L_{LT}, \]

donde el factor $\tfrac{1}{2}$ sobre la deuda de largo plazo recoge la convención empírica de KMV. Dados el valor de mercado del patrimonio $E$ (capitalización bursátil), su volatilidad $\sigma_E$, la tasa libre de riesgo $r$ y el horizonte $T$, se resuelve numéricamente el sistema de dos ecuaciones no lineales que proyecta el valor de los activos $\hat V_A$ y su volatilidad implícita $\hat\sigma_A$:

\[ E = \hat V_A\,\Phi(d_1) - D^{*}\,e^{-rT}\,\Phi(d_2), \qquad \sigma_E\,E = \Phi(d_1)\,\hat\sigma_A\,\hat V_A, \]

con $d_1 = \dfrac{\ln(\hat V_A/D^{*}) + \left(r + \tfrac{1}{2}\hat\sigma_A^{2}\right)T}{\hat\sigma_A\sqrt{T}}$ y $d_2 = d_1 - \hat\sigma_A\sqrt{T}$, donde $\Phi(\cdot)$ es la función de distribución normal estándar. La volatilidad del patrimonio se estima por \textbf{varianza realizada} \citep{BarndorffNielsenShephard2002}: para cada trimestre y banco, $\sigma_E = \sqrt{\sum_{t=1}^{Q} r_t^{2}}$, donde $r_t$ son los retornos logarítmicos diarios del patrimonio y $Q$ el número de días de transacción del trimestre. Con las soluciones $\hat V_A$ y $\hat\sigma_A$ se calcula la \textbf{distancia al incumplimiento} y la \textbf{frecuencia esperada de incumplimiento}:

\[ DD = \frac{\hat V_A - D^{*}}{\hat V_A\,\hat\sigma_A}, \qquad EDF = \Phi(-DD). \]

El valor $EDF$ de cada banco constituye su probabilidad incondicional de incumplimiento $p_i$, insumo del método de agregación.

\subsubsection{Factor sistémico y probabilidades condicionales}

La hipótesis identificadora de la agregación es la \textbf{independencia condicional}: los riesgos bancarios son independientes una vez condicionados a un único factor sistémico país, que se toma como el retorno del índice bursátil de cada economía. Sea $\rho_i$ la correlación entre los retornos del patrimonio del banco $i$ y los del factor sistémico $V$. Siguiendo la estructura unifactorial de \citet{Vasicek1997}, la probabilidad de incumplimiento \textbf{condicional} al estado $V$ del factor es

\[ p_i(V) = \Phi\!\left(\frac{\Phi^{-1}(p_i) - \rho_i\,V}{\sqrt{1-\rho_i^{2}}}\right). \]

Esta expresión traslada la dependencia de cola al canal sistémico: cuando $V$ adopta valores adversos, las probabilidades condicionales de todos los bancos aumentan simultáneamente, generando la correlación de pérdidas que un agregado de riesgos idiosincráticos ignoraría.

\subsubsection{Agregación de pérdidas por el método de punto de silla}

La pérdida agregada del sistema se define como $P = \sum_{i=1}^{N} e_i\,\mathbf{1}_{D_i}$, donde $e_i$ es la exposición al incumplimiento del banco $i$ (proxy: pasivos totales) y $\mathbf{1}_{D_i}$ la función indicadora de incumplimiento. El método de punto de silla \citep{MartinThompsonBrowne2001} traslada el problema al espacio de la función generadora de momentos, donde admite tratamiento analítico. La función generadora de cumulantes (CGF) de la pérdida, condicional al factor $V$, es

\[ K(s\mid V) = \sum_{i=1}^{N} \ln\!\left(1 - p_i(V) + p_i(V)\,e^{s\,a_i}\right), \]

con $a_i$ la exposición normalizada del banco $i$. Sus derivadas primera y segunda, $K'(s)$ y $K''(s)$, proporcionan la media y la varianza de la distribución de pérdidas. Para un nivel de pérdida $x$, el \textbf{punto de silla} $\hat t$ resuelve $K'(\hat t) = x$. La probabilidad de cola se aproxima entonces, siguiendo \citet{MartinThompsonBrowne2001}, por

\[ P(L > x \mid V) \approx 1 - \exp\!\left(K(\hat t\mid V) - \hat t\,x + \tfrac{1}{2}\hat t^{2} K''(\hat t)\right)\Phi\!\left(-|\hat t|\sqrt{K''(\hat t)}\right), \]

para $x$ por encima de la media (con la rama simétrica correspondiente en la cola inferior). La integración sobre la distribución del factor sistémico —supuesta $N(0,\sigma_V^{2})$— se resuelve por \textbf{cuadratura de Gauss–Hermite} con siete nodos:

\[ P(L > x) \approx \sum_{k=1}^{7} w_k\,P(L > x \mid V_k), \]

donde $\{V_k, w_k\}$ son los nodos y pesos de la cuadratura. La curva de probabilidad acumulada se evalúa sobre una grilla de 500 pasos entre las cotas de pérdida calibradas (1\% y 4,8\% de la exposición). El valor en riesgo al 99\% ($VaR_{99}$) se obtiene por interpolación lineal sobre esa curva, sin extrapolación fuera del rango de la grilla.

\subsubsection{Contribuciones marginales y definición de \texorpdfstring{$JLoss$}{JLoss}}

Siguiendo a \citet{MartinThompsonBrowne2001}, la contribución marginal del banco $n$ al riesgo del sistema, evaluada en el percentil 99, es

\[ C_n = e_n \cdot \frac{\sum_{k} \nu_k\,\tilde p_n(\hat t_k)}{\sum_{k}\nu_k}, \qquad \nu_k = w_k\,\frac{\exp\!\left(K(\hat t_k) - \hat t_k\,K'(\hat t_k)\right)}{\sqrt{K''(\hat t_k)}}, \qquad \tilde p_n(\hat t) = \frac{p_n\,e^{\hat t\,a_n}}{1 - p_n + p_n\,e^{\hat t\,a_n}}, \]

donde $\tilde p_n(\hat t)$ es la probabilidad exponencialmente inclinada (\textit{saddle-tilted}) del banco $n$. La suma de estas contribuciones en el percentil 99 constituye la \textbf{pérdida inesperada} ($UL$), mientras que la \textbf{pérdida esperada} es $EL = \sum_i e_i\,p_i$. La métrica de fragilidad se define finalmente como las pérdidas totales normalizadas por la exposición agregada:

\[ JLoss_{i,t} = \frac{EL_{i,t} + UL_{i,t}}{\sum_{j} EAD_{j,t}}\times 100, \]

donde la exposición al incumplimiento incorpora la pérdida dada el incumplimiento, $e_j = EAD_j\cdot LGD$, con $LGD = 45\%$ \citep{BIS2006}.

\subsubsection{Parametrización e implementación computacional}

El Cuadro 1 resume la parametrización, consistente con \citet{Chari2024b}: horizonte de un trimestre; $LGD = 45\%$; siete nodos de cuadratura; un único factor sistémico; cota inferior de pérdidas de 1,0\% y superior de 4,8\%; 500 pasos de discretización; percentil de 99\%; y factor sistémico $N(0,\,0{,}16\%)$. La totalidad del cómputo se reimplementó en Python validándolo contra el código original en MATLAB (\texttt{countrypd.m}, \texttt{loss\_distrib.m}, \texttt{loss\_contrib.m}, \texttt{get\_prob.m}, \texttt{LinealXY2.m}, \texttt{find\_saddlev1.m} y los módulos \texttt{get\_K*.m}), reproduciendo sus resultados. La búsqueda del punto de silla emplea el algoritmo de Brent (\texttt{brentq}) en lugar del Newton–Raphson original, por su mayor robustez ante curvaturas extremas de la CGF, e incorpora salvaguardas numéricas frente a desbordamiento en la evaluación de las exponenciales.

\vspace{1em}
\noindent \textbf{Cuadro 1 — Parámetros del método de punto de silla.} Horizonte: 1 trimestre $\cdot$ LGD: 45\% $\cdot$ Nodos de cuadratura: 7 $\cdot$ Factores sistémicos: 1 $\cdot$ Cota inferior de pérdidas: 0,010 $\cdot$ Cota superior: 0,048 $\cdot$ Pasos: 500 $\cdot$ Percentil: 0,99 $\cdot$ Factor sistémico: $N(0,\,0{,}16\%)$.
\vspace{1em}

\noindent\textbf{Una consecuencia de la calibración.} La malla de pérdidas está acotada
superiormente en el 4,8\% de la exposición y el valor en riesgo al 99\% se lee por
interpolación sobre esa malla, sin extrapolación. Un diagnóstico sobre la muestra revela que
esa cota es \textbf{activa en la práctica totalidad de los país-trimestre} (el 98\%): la
pérdida del percentil 99 del sistema casi siempre supera el 4,8\% de la exposición, de modo
que la componente de pérdida inesperada de $JLoss$ se evalúa en un nivel de estrés
\emph{fijo} y no en el percentil 99 verdadero. En consecuencia, (i) la variación de $JLoss$
—entre países y en el tiempo— está gobernada por la \textbf{pérdida esperada} (las
probabilidades de incumplimiento de Merton), no por la forma de la cola conjunta; y (ii) el
nivel de $JLoss$ está comprimido: recalculando la métrica con una malla más ancha, $JLoss$
escala por un factor $\approx 2{,}8$ (de forma heterogénea entre países), con una correlación
de $0{,}92$ con la serie base. Este trabajo reporta los resultados sobre la calibración
declarada por \citet{Chari2024b} y muestra en la Sección~6.5 que, al usar en su lugar la
serie de malla ancha, el coeficiente de interacción \textbf{conserva su magnitud y gana
precisión} (de $p=0{,}056$ a $p=0{,}032$): la censura de la cola, de tener algún efecto,
operaba en contra del hallazgo. En todo caso, $JLoss$ debe interpretarse de forma
\textbf{ordinal}.

\subsection{Medición del riesgo a la baja del crecimiento: la métrica \texorpdfstring{$GaR$}{GaR}}

$GaR_{i,t}$ es el cuantil de cola izquierda (5\%) de la distribución condicional del crecimiento del producto un trimestre adelante, estimado a partir de un índice de condiciones financieras mediante regresión cuantílica de panel. Su construcción adapta el marco de \textit{vulnerable growth} de \citet{ABG2019b} y la plataforma de CEMLA \citep{OssandonBusch2022}, con tres modificaciones metodológicas: estimación por programación lineal global, lectura directa del cuantil sin densidad kernel, y estimación pseudo–tiempo real.

\subsubsection{Índice de condiciones financieras (FCI)}

El FCI sintetiza por país un conjunto de indicadores de tensión financiera —de los mercados monetario, cambiario, accionario, de deuda soberana y bancario— junto con un bloque de tipo de cambio real efectivo. El procedimiento replica la cadena de rutinas del marco CEMLA y consta de tres pasos. Primero, cada indicador se \textbf{estandariza} mediante una transformación expansiva (\texttt{method\_b\_max}), consistente con la naturaleza pseudo–tiempo real de la estimación posterior. Segundo, se computan \textbf{correlaciones temporales por pares} entre los indicadores estandarizados mediante un suavizado exponencial (EWMA) con factor de decaimiento $\lambda = 0{,}85$. Tercero, el índice se agrega como una \textbf{forma cuadrática} sobre los indicadores estandarizados y la matriz de correlaciones EWMA $C_t$, siguiendo la lógica del indicador compuesto de estrés sistémico de \citet{HolloKremerLoDuca2012}:

\[ FCI_t = \sqrt{(w \circ s_t)'\,C_t\,(w \circ s_t)}, \]

donde $s_t$ es el vector de subíndices estandarizados y $w$ el vector de ponderaciones. La implementación se validó contra la salida de referencia de CEMLA (\texttt{FCI\_MEXICO.csv}), alcanzando una correlación de 0,9995 y un RMSE de 0,016 sobre 212 meses.

\subsubsection{Regresión cuantílica de panel con efectos fijos}

Para el cuantil $\tau$ y el horizonte $h$, la distribución condicional del crecimiento se caracteriza mediante la regresión cuantílica de panel

\[ Q\!\left(\Delta GDP_{i,t+h}\mid \cdot\,;\tau\right) = \alpha_h(\tau) + \beta_{1,h}(\tau)\,\Delta GDP_{i,t} + \beta_{2,h}(\tau)\,FCI_{i,t} + \beta_{3,h}(\tau)\,VIX_t + \mu_{i,h}(\tau), \]

con $h = 1$ trimestre y efectos fijos de país $\mu_i$ \citep{Koenker2004}. El VIX se ortogonaliza respecto de las condiciones locales. Los parámetros se obtienen minimizando la suma de pérdidas de comprobación,

\[ \min_{\beta,\,\mu}\;\sum_{i}\sum_{t}\rho_\tau\!\left(\Delta GDP_{i,t+h} - x_{i,t}'\beta - \mu_i\right), \qquad \rho_\tau(u) = u\,(\tau - \mathbf{1}\{u<0\}), \]

problema que se formula como un \textbf{único programa lineal} y se resuelve con el algoritmo HiGHS (vía \texttt{scipy.optimize.linprog}), superando la tolerancia del solver iterativo \texttt{rqpd} de R.

\subsubsection{De la función cuantílica a \texorpdfstring{$GaR$}{GaR}}

La regresión se estima sobre una grilla de $n_\tau = 39$ cuantiles, diseñada de modo que $\tau = 0{,}05$ caiga exactamente sobre un nodo. Para garantizar la monotonicidad se aplica el \textbf{reordenamiento} de \citet{ChernozhukovEtAl2010}:

\[ GaR_{i,t} \equiv Q\!\left(\Delta GDP_{i,t+1}\mid \cdot\,;\,\tau = 0{,}05\right). \]

\subsubsection{Estimación pseudo-tiempo real}

Para evitar el sesgo de anticipación (\textit{look-ahead bias}), la serie $\{GaR_{i,t}\}$ se construye con una \textbf{ventana expansiva}: en cada fecha $t$, los parámetros cuantílicos se reestiman utilizando exclusivamente la información disponible hasta $t$. La serie se estimó sobre un panel de 17 economías emergentes con historia amplia; la ventana estimable por país comienza tras los primeros trimestres de entrenamiento (Sección 5).

\subsubsection{Robustez: corroboración con densidad skew-\texorpdfstring{$t$}{t}}

Como control de robustez se ajusta la distribución \textit{skew-t} de \citet{ABG2019b} a los cuantiles condicionales estimados. La correlación entre esta serie y la obtenida por lectura directa del cuantil reordenado es de 0,9996.

\subsection{Estrategia de identificación econométrica}

\subsubsection{Especificación de referencia}

Siguiendo la lógica de identificación de \citet{Chari2024b}, la especificación de referencia es un modelo de panel con efectos fijos bidireccionales:

\[ CDS_{i,t} = \alpha_i + \gamma_t + \beta_1\,JLoss_{i,t} + \beta_2\,GaR_{i,t} + \omega'X_{i,t} + \epsilon_{i,t}, \]

donde $\alpha_i$ son efectos fijos de país que absorben toda heterogeneidad invariante en el tiempo; $\gamma_t$ son efectos fijos de tiempo que absorben los choques globales comunes —de modo que, bajo $\gamma_t$, los factores financieros globales quedan absorbidos y no se incluyen explícitamente en las especificaciones bidireccionales—; y $X_{i,t}$ es el vector de controles.

\subsubsection{Interacción y complementariedad}

La hipótesis central (H3) se contrasta incorporando el término de interacción:

\[ CDS_{i,t} = \alpha_i + \gamma_t + \beta_1\,JLoss_{i,t} + \beta_2\,GaR_{i,t} + \theta\,\bigl(JLoss_{i,t}\times GaR_{i,t}\bigr) + \omega'X_{i,t} + \epsilon_{i,t}. \]

Bajo la convención de signo adoptada, la complementariedad supra-aditiva predice $\theta < 0$. El efecto marginal de la fragilidad sobre el diferencial,

\[ \frac{\partial\,CDS_{i,t}}{\partial\,JLoss_{i,t}} = \beta_1 + \theta\,GaR_{i,t}, \]

se reporta evaluado a distintos niveles de $GaR$ en un gráfico de efectos parciales análogo al de \citet{Chari2024b}.

\subsubsection{Modelo de umbral de panel}

La versión no lineal de la hipótesis de amplificación se contrasta mediante un modelo de umbral \citep{Hansen1999,Hansen2000}:

\[ CDS_{i,t} = \alpha_i + \gamma_t + \boldsymbol{\beta}^{(1)'} Z_{i,t}\,\mathbf{1}\{q_{i,t}\le\lambda\} + \boldsymbol{\beta}^{(2)'} Z_{i,t}\,\mathbf{1}\{q_{i,t}>\lambda\} + \omega'X_{i,t} + \epsilon_{i,t}, \]

donde $q_{i,t}=GaR_{i,t}$. El umbral $\lambda$ se estima por mínimos cuadrados concentrados tras la transformación \textit{within}; su significancia se evalúa mediante un contraste de razón de verosimilitud cuya distribución se aproxima por \textit{bootstrap} \citep{Hansen2000}.

\subsubsection{Construcción del panel e inferencia}

Las series de $JLoss$, $GaR$, CDS soberano y controles se consolidan por país y trimestre. La inferencia se basa en errores estándar de Driscoll–Kraay \citep{DriscollKraay1998} ($vcovSCC$), robustos simultáneamente a heterocedasticidad, autocorrelación serial y dependencia transversal —justificados por los diagnósticos de la Sección 6.6—, contrastados frente a \textit{clustering} por país, por tiempo y al estimador HAC de Arellano.

\section{Datos}

\subsection{Un solo panel}

Este trabajo descansa sobre \textbf{un único panel} de economías emergentes a frecuencia trimestral, sin partición en subconjuntos. El principio de construcción es homogeneidad de fuente: todos los insumos micro y de mercado —balances bancarios, capitalización bursátil, CDS soberano, VIX y factores financieros globales— provienen de la terminal Bloomberg, extraídos con un mismo protocolo para las veinte economías del universo objetivo; solo los componentes reales del índice de condiciones financieras que alimenta $GaR$ (IPC, PIB, índice accionario, tipo de cambio real efectivo, rendimiento soberano a diez años) se mantienen de estadísticas nacionales, dado que Bloomberg no es fuente primaria de cuentas nacionales. El Anexo~\ref{chap:anexoB} tabula la procedencia de cada serie.

El \textbf{criterio de inclusión} es deliberadamente permisivo: una economía entra al panel si se dispone de $JLoss$; su spread soberano se mide \textbf{siempre y solo} por el CDS soberano a cinco años de Bloomberg, y allí donde ese dato no existe la celda queda vacía —no se sustituye por el EMBI de bonos ni por ningún \textit{proxy} de spread, para no mezclar instrumentos heterogéneos en la variable dependiente. Una economía sin CDS, o sin $GaR$, permanece en el panel y en la figura de cobertura pero no aporta observaciones a la estimación del mecanismo. La \textbf{única exclusión} es Corea del Sur: el valor de mercado de su banca cotizada es persistentemente cercano al 4\,\% del punto de incumplimiento —frente a un rango de 15--34\,\% en el resto del panel—, reflejo del descuento estructural de valoración de la banca coreana (\textit{Korea discount}) y no de su solvencia; el modelo de Merton lee ese descuento como incumplimiento inminente y devuelve un $JLoss$ implausible (mediana 25, creciente a 47). Es una exclusión por insumo no válido, no por dato faltante, y coincide con la reconstrucción previa basada en datos regulatorios, que tampoco incluía a Corea.

Junto a Corea, se excluye también Bulgaria por la misma razón de fondo —$JLoss$ que no es una medida a nivel de país—: solo un banco búlgaro cotiza (FIBank), de modo que la métrica se construye sobre una única institución (los 76 trimestres quedan bajo el mínimo de bancos requerido) y arroja una mediana de 29, incoherente con un sistema bancario solvente. La muestra de estimación resultante reúne las economías con CDS soberano, $JLoss$ y $GaR$ disponibles simultáneamente: \textbf{catorce países y 838 observaciones país-trimestre} (2004Q1--2026Q1). Once de ellos —Brasil, Chile, China, Colombia, Filipinas, Indonesia, Malasia, México, Perú, Sudáfrica y Turquía— tienen CDS de historia continua y contribuyen entre 40 y 89 trimestres cada uno (el mínimo, Malasia, limitado por la ventana de entrenamiento del $GaR$); Hungría y Polonia solo tienen precios de CDS en la ventana 2012Q3--2015Q4 (14 trimestres cada uno) y Pakistán apenas uno reciente. Argentina, Egipto y Rusia tienen $JLoss$ y CDS pero no $GaR$; India tiene $JLoss$ y $GaR$ pero no emite CDS soberano de referencia. La Figura~\ref{fig:cobertura} traza esta cobertura.

\begin{figure}[t]
\centering
\includegraphics[width=\textwidth]{fig_cobertura.pdf}
\caption[Cobertura del panel, por país]{Cobertura del panel, por país. Marca gris: trimestres con $JLoss$. Marca celeste: además con CDS soberano de Bloomberg. Marca azul: muestra de estimación del mecanismo (CDS + $JLoss$ + $GaR$ simultáneamente). Los países ordenados de menor a mayor aporte a la estimación.}
\label{fig:cobertura}
\end{figure}

\subsection{Variable dependiente}
El spread soberano se mide \emph{exclusivamente} por el CDS soberano a cinco años en dólares (Bloomberg; \textit{ticker} \texttt{<emisor> CDS USD SR 5Y Corp}, con la curva Markit \texttt{CCHIN1U5} para China), en puntos básicos y promediado a frecuencia trimestral. El CDS, a diferencia del EMBI de bonos, aísla el componente puro de riesgo de crédito y es homogéneo entre países. Donde no hay CDS, la observación no entra a la regresión.

\subsection{Insumos de \texorpdfstring{$JLoss$}{JLoss}}
Los insumos micro son las hojas de balance bancarias trimestrales (activos totales, patrimonio, deuda emitida, ingreso neto, ingresos, caja) y la capitalización bursátil diaria de los bancos cotizados, de Bloomberg, con un protocolo de extracción común a los 113 bancos. La volatilidad del patrimonio se calcula por varianza realizada de los retornos diarios; la cotización utilizada es siempre la de la bolsa local, nunca la del ADR, y se verifica contra la moneda del balance. El motor de agregación es \texttt{jloss\_engine.py}; los campos Bloomberg exactos figuran en el Anexo~\ref{chap:anexoB}.

\subsection{Insumos de \texorpdfstring{$GaR$}{GaR}}
Los componentes del FCI por país (IPC, rendimiento soberano a diez años, índice accionario, REER) provienen de estadísticas nacionales; el VIX de Bloomberg entra como \textit{proxy} de condiciones financieras globales, ortogonalizado respecto de las condiciones locales.

\subsection{Variables de control}
Los controles domésticos son la deuda del gobierno general/PIB y el balance fiscal/PIB (FMI, \textit{World Economic Outlook}), las reservas/PIB y la cuenta corriente/PIB (Banco Mundial), la inflación interanual (de la serie de IPC) y el tipo de cambio real efectivo. Se reconstruyeron para \textbf{todos} los países del panel, de modo que la especificación con controles se estima sobre la muestra completa. Los factores globales son el VIX, la tasa del Tesoro estadounidense a diez años y el spread de alto rendimiento estadounidense (Bloomberg).

\subsection{Estadística descriptiva}

El Cuadro 2 resume las variables centrales sobre la muestra de estimación. Un rasgo relevante para lo que sigue: a diferencia de la reconstrucción con datos regulatorios —en la que Brasil concentraba el grueso de la fragilidad del panel ($JLoss$ medio $\approx$16)— la serie construida homogéneamente desde Bloomberg es \textbf{mucho más comprimida en el corte transversal} (desviación estándar $\approx$4 frente a 8,3), con Turquía y China, no Brasil, en el extremo superior. Esta homogeneidad favorece la comparabilidad entre países pero, como se discute en el Capítulo~\ref{chap:paper1}, reduce la variación transversal de la fragilidad que identifica los términos de interacción de orden superior.

\begin{center}
\footnotesize
\begin{tabular}{lrr}
\toprule
Variable & Media & Desv. \\
\midrule
CDS soberano 5A (pb) & 149 & 102 \\
JLoss                & 4,1  & 4,1  \\
GaR (q05)            & 0,010 & 0,040 \\
Deuda gob.\ gral./PIB (\%) & 46 & 20 \\
Bal.\ fiscal/PIB (\%) & $-3,0$ & 2,3 \\
Reservas/PIB (\%)     & 19 & 9 \\
C.\ corriente/PIB (\%) & $-1,1$ & 3,0 \\
Inflación YoY (\%)    & 5,6 & 8,4 \\
REER (índice, 2010=100) & 108 & 20 \\
\bottomrule
\end{tabular}

\medskip
\footnotesize\textbf{Cuadro 2.} Estadísticos descriptivos sobre la muestra de estimación (observaciones con CDS + $JLoss$ + $GaR$). Fuente y transformación de cada serie: Anexo~\ref{chap:anexoB}.
\end{center}

\begin{figure}[t]
\centering
\includegraphics[width=\textwidth]{fig_jloss_paises.pdf}
\caption[Serie de JLoss por país del roster]{Serie trimestral de $JLoss$ por país del \textit{roster}, panel Bloomberg (motor de punto de silla, $\rho=0{,}4$). La escala vertical es común salvo cuando la serie del país la excede.}
\label{fig:jlosspaises}
\end{figure}

\section{Resultados}

\subsection{Estrategia de identificación}

La estimación se realiza sobre la muestra descrita en la Sección 5: las observaciones del panel único con CDS soberano, $JLoss$ y $GaR$ disponibles simultáneamente ($N=838$, catorce países; $N=738$ y trece países bajo la especificación con controles, que pierde el único trimestre de Pakistán). La especificación de referencia incorpora efectos fijos bidireccionales —de país y de tiempo—, de modo que los efectos de país absorben la heterogeneidad estructural no observada y los efectos de tiempo capturan la totalidad de los choques globales comunes. La inferencia se basa en errores estándar de Driscoll–Kraay, contrastada con \textit{clustering} por país y por tiempo y con \textit{wild cluster bootstrap}. La hipótesis central —complementariedad entre fragilidad bancaria sistémica y riesgo de cola izquierda del crecimiento— se contrasta a través del término de interacción $\theta$ en

\[ CDS_{i,t} = \alpha_i + \gamma_t + \beta_1\,JLoss_{i,t} + \beta_2\,GaR_{i,t} + \theta\,\bigl(JLoss_{i,t}\times GaR_{i,t}\bigr) + \omega'X_{i,t} + \epsilon_{i,t}, \]

donde $X_{i,t}$ es el vector de seis controles domésticos. Bajo la convención de signos adoptada, la predicción teórica del mecanismo de \textit{doom loop} es inequívoca: $\theta<0$.

\subsection{Resultado principal: el canal de complementariedad}

\begin{table}[t]
\centering
\small
\setlength{\tabcolsep}{4.5pt}
\begin{tabular}{lrrrrr}
\toprule
 & M1 & \textbf{M2} & M3 & cluster país & cluster tiempo \\
 & (sin contr.) & \textbf{(+contr.)} & (+glob.) & & \\
\midrule
$JLoss\times GaR$ ($\theta$) & $-0,543^{**}$ & $\mathbf{-0,354^{*}}$ & $-0,339$ & $-0,354^{***}$ & $-0,354^{**}$ \\
$t$ (Driscoll–Kraay)         & $-2,20$ & $\mathbf{-1,91}$ & $-1,65$ & $-3,39$ & $-2,27$ \\
$p$                          & $0,028$ & $\mathbf{0,056}$ & $0,100$ & $0,001$ & $0,023$ \\
$p$ (\textit{wild bootstrap})& $0,127$ & $\mathbf{0,035}$ & $0,039$ & --- & --- \\
$N$                          & 838 & 738 & 738 & 738 & 738 \\
$R^2_{\text{within}}$        & $0,15$ & $0,18$ & $0,63$ & $0,18$ & $0,18$ \\
\bottomrule
\end{tabular}
\smallskip
\small\textbf{Cuadro 3.} Coeficiente de interacción $\theta$ sobre el panel único (14 países; M2, la especificación de referencia, usa 13 al perder Pakistán, que aporta un solo trimestre efectivo). M1: efectos fijos país+tiempo sin controles. M2: + seis controles domésticos. M3: efectos fijos de país + factores financieros globales explícitos (sin efectos de tiempo). Errores de Driscoll–Kraay salvo las dos últimas columnas; *** $p<0{,}01$, ** $p<0{,}05$, * $p<0{,}10$. Ver \texttt{1\_Codigo/Panel/bbg/NUMEROS\_CANONICOS\_BBG.md}.
\end{table}

El Cuadro 3 presenta el resultado central. En la especificación de referencia M2 —con el vector completo de controles macroeconómicos domésticos— el coeficiente de interacción $\hat\theta=-0{,}354$ es negativo, el signo predicho por el multiplicador de \textit{doom loop}. Su significancia es \textbf{marginal} bajo errores de Driscoll–Kraay ($t=-1{,}91$, $p=0{,}056$) y bajo \textit{wild cluster bootstrap} ($p=0{,}035$), y \textbf{convencional} bajo \textit{clustering} por país ($p=0{,}001$) y por tiempo ($p=0{,}023$). El signo es robusto a la especificación: sin controles (M1) la magnitud es mayor ($\hat\theta=-0{,}543$, $p=0{,}028$); con los factores globales explícitos en lugar de los efectos de tiempo (M3), $\hat\theta=-0{,}339$ ($p\approx0{,}10$, justo en el límite convencional). El ajuste dentro del panel es adecuado ($R^2_{\text{within}}=0{,}18$).

Esta es una evidencia más modesta que la de una reconstrucción previa con datos regulatorios (donde $\hat\theta\approx-0{,}34$ era significativo al 5\% de forma robusta), y se reporta como tal: lo que el panel homogéneo respalda con solidez es el \textbf{signo} y la \textbf{forma} de la complementariedad —confirmada por el efecto marginal (Sección 6.3) y por el modelo de umbral (Sección 6.4)—, no una magnitud puntual con significancia a prueba de toda inferencia.

\subsection{Magnitud económica: el multiplicador del riesgo soberano}

La relevancia económica del mecanismo se aprecia en el \textbf{efecto marginal} de la fragilidad bancaria condicionado al estado del riesgo de cola,
\[\frac{\partial\,CDS_{i,t}}{\partial\,JLoss_{i,t}}=\hat\beta_1+\hat\theta\,(GaR_{i,t}-\overline{GaR}),\]
que la Figura~\ref{fig:efmarg} traza. Evaluado a partir de M2, una unidad adicional de $JLoss$ eleva el CDS soberano en $\approx +4{,}6$ puntos básicos cuando el riesgo de cola está en su percentil 10 (cola severa), frente a $\approx +1{,}8$ en el percentil 90 (entorno benigno). El efecto es monótonamente decreciente en $GaR$ —la firma empírica de la amplificación supra-aditiva— y la banda de confianza al 90\% excluye el cero en todo el tramo de cola severa.

\begin{figure}[t]
\centering
\includegraphics[width=0.8\textwidth]{fig_efecto_marginal.pdf}
\caption[Efecto marginal de la fragilidad según el riesgo de cola]{Efecto marginal de la fragilidad bancaria sobre el CDS soberano según el riesgo de cola del crecimiento (M2). Banda: IC 90\% (errores Driscoll–Kraay). Eje horizontal invertido: cola más severa a la derecha.}
\label{fig:efmarg}
\end{figure}

\subsection{No linealidad: el modelo de umbral de Hansen}

Como contraste no lineal alternativo se estima un modelo de umbral de panel \citep{Hansen1999,Hansen2000}, con el propio $GaR$ como variable de umbral y sin imponer la forma funcional de la interacción. El umbral estimado es $\hat\gamma\approx -0{,}15$ puntos porcentuales de $GaR$: en el régimen de cola severa ($GaR\le\hat\gamma$) el efecto de la fragilidad sobre el spread es de $+8{,}1$ puntos básicos por unidad de $JLoss$, frente a $+2{,}3$ en el régimen benigno (Figura~\ref{fig:umbral}); el contraste de razón de verosimilitud es amplio ($LR=80$). El efecto de la fragilidad bancaria sobre el riesgo soberano es, por tanto, unas \textbf{tres veces y media mayor} cuando el crecimiento ya está en su cola izquierda: el modelo de umbral corrobora el hallazgo del término multiplicativo.

\begin{figure}[t]
\centering
\includegraphics[width=0.55\textwidth]{fig_umbral.pdf}
\caption[Modelo de umbral de Hansen]{Modelo de umbral de panel (Hansen): efecto de $JLoss$ sobre el CDS soberano por régimen de $GaR$.}
\label{fig:umbral}
\end{figure}

\subsection{Robustez}

\subsubsection{Invariancia a la medida de cola}

El coeficiente de interacción es esencialmente idéntico bajo tres construcciones alternativas del riesgo de cola: $\hat\theta=-0{,}354$ leyendo $GaR$ directamente de la función cuantil ($p=0{,}056$); $\hat\theta=-0{,}354$ ajustando la densidad \textit{skew-t} de \citet{ABG2019b} a los cuantiles condicionales ($p=0{,}051$); y $\hat\theta=-0{,}390$ sustituyendo el cuantil 5\% por el \textit{Expected Shortfall} ($p=0{,}036$). El resultado \textbf{no depende de la especificación de la primera etapa} con la que se estima el riesgo de cola.

\subsubsection{Cuándo se identifica la interacción: una regularidad post-GFC}

Estimando $\theta$ sobre ventanas móviles de cinco años se observa un patrón nítido (Figura~\ref{fig:ventanas}): la interacción es \textbf{positiva o nula} en las ventanas que cubren la crisis financiera global (2006--2010: $\hat\theta=+1{,}1$, $t=0{,}8$; 2009--2013: $\hat\theta=+0{,}3$, $t=0{,}5$) y \textbf{negativa y significativa en todas las ventanas que empiezan en 2012 o después} (2012--2016: $\hat\theta=-1{,}2$, $t=-3{,}3$; 2015--2019: $-0{,}4$, $t=-1{,}6$; 2018--2022: $-0{,}2$, $t=-2{,}2$; 2021--2025: $-0{,}3$, $t=-3{,}1$). La complementariedad es, por tanto, un fenómeno del período \textbf{posterior a la crisis financiera global}, estable desde 2012 hasta el final de la muestra.

\begin{figure}[t]
\centering
\includegraphics[width=0.85\textwidth]{fig_ventanas_theta.pdf}
\caption[Coeficiente de interacción por ventana móvil]{Coeficiente de interacción $\theta$ ($JLoss\times GaR$, especificación M2) estimado sobre ventanas móviles de cinco años. Barras rojas: $|t|<1{,}64$. El signo negativo aparece —y se vuelve significativo— en todas las ventanas que empiezan en 2012 o después.}
\label{fig:ventanas}
\end{figure}

Esto matiza —y corrige— una lectura tentadora. Restringida al período \textbf{previo a 2020}, la interacción conserva el signo pero pierde precisión ($\hat\theta=-0{,}39$, $t=-1{,}02$); ello \emph{no} se debe a que el mecanismo sea un artefacto de la pandemia, sino a que la submuestra pre-2020 promedia la señal post-2012 con el período \textbf{2004--2011}, en el que la interacción era nula o de signo contrario. La prueba directa lo confirma: al interactuar $JLoss\times GaR$ con una variable indicadora del período posterior a 2020, el término adicional es $+0{,}46$ ($t=0{,}8$, no significativo) —no hay un quiebre discreto en 2020. La lectura correcta es que la covariación entre fragilidad bancaria, riesgo de cola y spread soberano que identifica el mecanismo es \textbf{tenue en el período de expansión previo a la crisis y nítida desde entonces}, coherente con un multiplicador que solo opera cuando el sistema financiero global ha dejado atrás su fase de baja percepción de riesgo.

\subsubsection{El corte transversal efectivo: la identificación descansa en pocos países}

Reestimando $\theta$ excluyendo secuencialmente cada país (\textit{leave-one-country-out}, Figura~\ref{fig:forest}), el coeficiente permanece negativo en las catorce submuestras, y negativo en las noventa y un submuestras que excluyen \emph{dos} países a la vez. Pero su \textbf{significancia} descansa de forma desproporcionada en un país: al excluir China, $\hat\theta$ cae a $-0{,}17$ ($t=-1{,}33$, no significativo), y al excluir China y Turquía juntas la interacción se anula ($\hat\theta=-0{,}07$, $t=-0{,}3$). El diagnóstico por país explica por qué (Anexo~\ref{chap:anexoB}): nueve de los trece países del panel de estimación tienen un $JLoss$ mediano entre 2,2 y 2,6 —casi sin variación transversal—, mientras que China (mediano 3,8; máximo 29), Turquía (mediano 6,0; máximo 28) y Brasil (mediano 4,0) se separan claramente de ese conglomerado y aportan la mayor parte de la dispersión de la fragilidad entre países. El término de interacción se identifica, en la práctica, de esos tres o cuatro países con movimiento sustantivo de $JLoss$ y de los episodios de estrés macrofinanciero: la amplitud nominal de "catorce países" sobreestima el corte transversal que efectivamente identifica el mecanismo.

\subsubsection{Otras pruebas}

Excluyendo la deuda/PIB del vector de controles la interacción se mantiene ($\hat\theta=-0{,}33$, $t=-2{,}49$). Un modelo dinámico —añadiendo el CDS rezagado— es inapropiado para la forma de este panel ($T\approx88$, $N=13$: el sesgo de Nickell es del orden de $1/T\approx1\%$, y los estimadores GMM de panel dinámico están diseñados para el caso opuesto, $N$ grande y $T$ corto), pero devuelve el mismo signo ($\hat\theta_{JxG}=-0{,}25$). Un \textit{placebo} que destruye por completo la estructura temporal y transversal de $GaR$ mediante reasignación aleatoria deja el $\hat\theta$ observado en el percentil~5 de la distribución nula (Sección~6.7). Recalculando $JLoss$ con una malla de pérdidas más ancha —que corrige la censura de la cota del 4,8\% discutida en la Sección~4.1—, la métrica escala por un factor $\approx 2{,}8$ de forma heterogénea entre países (correlación $0{,}92$ con la serie base). El coeficiente de interacción sobre esta serie corregida es $\hat\theta=-0{,}330$ ($t=-2{,}15$, $p=0{,}032$) en M2 y $\hat\theta=-0{,}572$ ($t=-2{,}57$, $p=0{,}010$) en M1: prácticamente la misma magnitud que en la especificación base y con \textbf{mayor} precisión. La censura de la cola, de haber sesgado algo, sesgaba en contra del hallazgo.

\begin{figure}[t]
\centering
\includegraphics[width=0.72\textwidth]{fig_forest_theta.pdf}
\caption[Coeficiente de interacción theta: especificaciones y robustez]{Coeficiente de interacción $\theta$ por especificación y prueba de robustez; barras: IC 95\% (errores de Driscoll--Kraay). El punto estimado es negativo en las veinticuatro filas —todas las especificaciones, todas las submuestras temporales y todos los \textit{leave-one-country-out}—; la submuestra previa a 2020 conserva el signo pero su intervalo cruza el cero.}
\label{fig:forest}
\end{figure}

\subsection{Controles macroeconómicos y diagnósticos del panel}

Los controles de fondos y de fundamentos entran, en su mayoría, con los signos esperados: la deuda del gobierno general sobre PIB es positiva y significativa ($t=4{,}5$), la inflación interanual positiva y significativa ($t=2{,}7$), y el tipo de cambio real efectivo negativo y fuertemente significativo ($t=-6{,}8$) —la apreciación real, asociada a entradas de capital y fortaleza de fundamentos, reduce la prima de riesgo. El balance fiscal y las reservas entran con el signo esperado pero no son significativos. La cuenta corriente es la única con signo contraintuitivo (mejor saldo asociado a mayor spread), y tampoco es significativa; un patrón plausiblemente atribuible a \textbf{endogeneidad por causalidad inversa}. Que la interacción $JLoss\times GaR$ conserve el signo predicho una vez condicionada por este bloque —incluida una deuda/PIB que aquí sí pesa— es evidencia de que la complementariedad no es un artefacto de variables omitidas de naturaleza fiscal o externa.

La prueba de Pesaran \citep{Pesaran2021} de dependencia transversal no rechaza la independencia de los residuos ($CD=-1{,}2$, $p=0{,}23$) una vez incluidos los efectos fijos de tiempo; se mantienen de todos modos los errores de Driscoll–Kraay por la autocorrelación serial detectada (correlación de primer orden de los residuos $\approx 0{,}8$). La prueba de Hausman rechaza la consistencia del estimador de efectos aleatorios. Los factores de inflación de varianza son uniformemente bajos.

\subsection{Identificación causal: qué se puede y qué no se puede afirmar}

Los resultados anteriores son, en rigor, asociaciones condicionales robustas en signo bajo un conjunto exigente de controles y efectos fijos, no estimaciones con identificación causal cerrada. Tres estrategias adicionales precisan qué margen de esa distinción es defendible.

\textbf{Wild cluster bootstrap} \citep{CameronGelbachMiller2008}. Con trece \textit{clusters} de país, la inferencia asintótica sobre errores agrupados puede ser optimista. El $p$-valor de $\theta$ (M2) bajo \textit{wild cluster bootstrap} (999 repeticiones) es $0{,}035$: marginalmente significativo, más ajustado que bajo Driscoll–Kraay pero por debajo del umbral convencional del 5\% por muy poco margen.

\textbf{Variables instrumentales por \textit{shift-share}} \citep{GoldsmithPinkhamSorkinSwift2020}. Instrumentando el nivel de $JLoss$ con el diferencial de liquidez de fondeo del Tesoro estadounidense (\textit{on/off-run}) interactuado con la exposición pre-muestral de cada país —estimada en el período previo a 2012, el mismo quiebre de régimen documentado en la Sección~6.5—, la primera etapa es \textbf{fuerte} ($F=21{,}6$) y el efecto de nivel instrumentado es positivo y muy significativo ($\hat\beta=+17{,}5$ pb por unidad de $JLoss$, $p=0{,}001$), de la dirección causal esperada (banco$\to$soberano). Un segundo choque global —el tipo de cambio efectivo nominal amplio del dólar (BIS, canasta de 64 economías), el instrumento estándar del canal de descalce cambiario— tiene una primera etapa igualmente fuerte ($F=23{,}7$) pero un efecto de segunda etapa pequeño y no significativo ($p=0{,}46$) cuando se usa solo; combinando ambos en una especificación sobre-identificada, el test de Sargan \textbf{rechaza} la validez conjunta de los dos instrumentos ($p=0{,}003$). La lectura honesta es que el diferencial de liquidez de fondeo bancario —con una restricción de exclusión más estrecha, al operar específicamente sobre el costo de fondeo de los bancos— es el instrumento más creíble de los dos, y que el rechazo de Sargan advierte contra tratar el efecto de nivel como definitivamente identificado; con ese instrumento solo, la evidencia de la dirección causal banco$\to$soberano es \textbf{sólida}.

\textbf{Regresores generados.} Tanto $JLoss$ (modelo estructural de \citet{Merton1974b} más punto de silla) como $GaR$ (regresión cuantílica de panel) son variables \emph{estimadas}, no observadas, de modo que la inferencia de segunda etapa que las trata como datos subestima la varianza de $\hat\theta$. Dos consideraciones acotan el problema. Primera, el error de medición clásico en un regresor \textbf{atenúa} el coeficiente hacia cero: el $\hat\theta$ reportado es un límite conservador y el efecto verdadero es, si acaso, de mayor magnitud. Segunda, perturbando la serie de $GaR$ con ruido gaussiano de desviación estándar creciente, $\hat\theta$ es estable hasta perturbaciones del orden del 25\,\% de la desviación de $GaR$ ($\hat\theta=-0{,}34$) y solo se atenúa a $-0{,}29$ con una perturbación del 50\,\%, sin cambiar de signo en ninguna réplica. Un \textit{bootstrap} que reestime la regresión cuantílica de $GaR$ en cada réplica —computacionalmente intensivo por la ventana expansiva— queda como refinamiento pendiente.

\textbf{Placebo.} Reasignando aleatoriamente la serie completa de $GaR$ entre observaciones —lo que destruye toda su estructura y genera la distribución nula exacta de $\hat\theta$— el coeficiente observado ($-0{,}354$) cae en el percentil~5 de esa distribución (media nula $\approx0$, desviación $0{,}20$): un $p$ de una cola $\approx0{,}05$, coherente con la inferencia analítica. Un placebo que solo permuta $GaR$ \emph{dentro} de cada país deja una distribución nula centrada en $-0{,}27$, no en cero: parte de la identificación proviene de la covariación \emph{entre} países de la fragilidad y el riesgo de cola —que los efectos fijos bidireccionales no absorben del todo— y no solo de la sincronía intra-país, una matización que conviene explicitar.

\textbf{Proyecciones locales} \citep{Jorda2005} y \textbf{triple interacción institucional}. La respuesta dinámica del spread a un choque de fragilidad, condicionada al régimen de $GaR$, es ruidosa y no confirma una amplificación dinámica creciente en el horizonte. La triple interacción $JLoss\times GaR\times$Institución (calidad de gobierno, WGI) no es significativa ($t=-0{,}10$); los datos institucionales empleados son además provisionales.

\textbf{Síntesis.} Lo defendible con los datos y estrategias disponibles hoy es: (i) el mecanismo de complementariedad es teóricamente sólido y su signo y forma son visibles a través de especificaciones, del efecto marginal, del modelo de umbral y de tres construcciones alternativas del riesgo de cola; (ii) la dirección causal de nivel banco$\to$soberano cuenta con evidencia sólida bajo un instrumento fuerte y de restricción de exclusión creíble, con la salvedad de que un segundo instrumento, también fuerte en la primera etapa, no corrobora la magnitud y falla el test de sobre-identificación al combinarse con el primero; (iii) la significancia \textit{puntual} de la interacción $\theta$ es marginal —$p\approx0{,}05$ bajo Driscoll--Kraay, bajo \textit{wild cluster bootstrap} y bajo el placebo de reasignación— y su identificación descansa en el período posterior a la crisis financiera global y en el puñado de países con variación sustantiva de $JLoss$. La contribución empírica es, por tanto, el signo y la coherencia del mecanismo, no una magnitud puntual robusta a toda prueba, en una muestra —economías emergentes con CDS soberano y banca cotizada— cuya profundidad transversal efectiva y cuya profundidad temporal en tiempos tranquilos siguen siendo el cuello de botella.

\subsection{Síntesis}

Sobre un único panel de catorce economías emergentes con datos homogéneos de Bloomberg, la evidencia respalda el \textbf{signo} y la \textbf{forma} de la hipótesis de complementariedad: el coeficiente de interacción $JLoss\times GaR$ sobre el CDS soberano es negativo en la especificación de referencia y en todas sus variantes, negativo en las catorce submuestras \textit{leave-one-out}, con un efecto marginal de la fragilidad monótonamente creciente a medida que el riesgo de cola se agrava y un modelo de umbral que corrobora la no linealidad sin imponerla. La \textbf{significancia puntual} de $\theta$ es marginal —$p=0{,}056$ bajo errores de Driscoll–Kraay, $p=0{,}035$ bajo \textit{wild cluster bootstrap}, $p=0{,}001$ bajo \textit{clustering} por país— y su identificación tiene dos fronteras: es un fenómeno \textbf{posterior a la crisis financiera global} —negativo y significativo en todas las ventanas móviles que empiezan en 2012, nulo o positivo antes— y descansa de forma desproporcionada en los pocos países con variación sustantiva de la fragilidad (al excluir China la interacción deja de ser significativa). La batería de identificación causal sugiere una dirección de nivel banco$\to$soberano sin cerrarla. Esta estructura de resultados —signo y forma robustos, magnitud puntual marginalmente significativa, identificación acotada en el tiempo y en el corte transversal— es el resumen más honesto que la evidencia de este panel permite ofrecer.

\section{Discusión y conclusiones del capítulo}

\subsection{Interpretación económica}

El hallazgo central de este trabajo —un coeficiente de interacción $JLoss\times GaR$ negativo, del signo predicho, con significancia marginal y corroborado por un modelo de umbral que no impone su forma funcional— admite una interpretación económica directa en términos del mecanismo de \textit{doom loop}. La fragilidad del sistema bancario opera como una aproximación al pasivo contingente que el soberano podría verse forzado a asumir mediante un rescate; pero ese pasivo sólo se torna relevante para la prima de riesgo cuando la economía transita un estado de cola. Que el efecto de la fragilidad bancaria sobre el spread sea unas tres veces y media mayor en el régimen de cola severa que en el benigno es la firma empírica del multiplicador: no la suma de dos vulnerabilidades, sino su producto. Que la interacción emerja como regularidad solo en el período posterior a la crisis financiera global —y no en los años de expansión previos— es consistente con esa misma lectura: el multiplicador opera cuando el precio del riesgo soberano es sensible a la cola, no en fases de baja percepción de riesgo global.

\subsection{Implicancias de política}

Si la transmisión de la fragilidad bancaria al costo de financiamiento soberano se concentra en los estados de cola del crecimiento, la regulación orientada a la resiliencia del sistema bancario —capital, liquidez, provisiones contracíclicas— reduce la exposición del soberano a los escenarios macroeconómicos adversos y su prima de riesgo en los mercados internacionales. La fragilidad bancaria y el riesgo de cola del crecimiento deben monitorearse conjuntamente: la coincidencia de ambos, y no cada uno por separado, es la señal de alarma relevante. La evidencia de este panel establece esa implicancia a nivel de \textbf{signo y mecanismo}; su magnitud precisa —cuánto exactamente amplifica cada unidad de fragilidad el spread bajo estrés— exige la profundidad temporal en tiempos normales que el panel todavía no tiene.

\subsection{Limitaciones}

Primero, la interacción se identifica en el período \textbf{posterior a la crisis financiera global}: negativa y significativa en todas las ventanas móviles de cinco años que empiezan en 2012, pero nula o de signo contrario en 2004--2011, de modo que el estimador de muestra completa previa a 2020 resulta impreciso. El panel tiene amplitud nominal (catorce países) y una ventana hasta 2026, pero pocos episodios de cola independientes. Segundo, y relacionado: nueve de los trece países del panel de estimación casi no exhiben variación de $JLoss$ (mediana 2,2--2,6, máximo por debajo de 18), de modo que el término de interacción se identifica en la práctica de China, Turquía y Brasil más los episodios de estrés; al excluir China la significancia se pierde. La homogeneización de la métrica de fragilidad al construirla desde una única fuente —ventaja para la comparabilidad— comprime esa dispersión transversal (desviación estándar $\approx$4 frente a 8,3 con datos regulatorios) y es, como se discute en el Capítulo~\ref{chap:paper1}, la razón más probable de que la prueba de amplificación por concentración quede sin potencia. A ello se suma la calibración del propio motor: la cota superior de la malla de pérdidas (4,8\% de la exposición) es activa en el 98\% de las observaciones, de modo que $JLoss$ mide en la práctica la pérdida esperada más una contribución de cola evaluada a un nivel de estrés fijo; recalculada con una malla más ancha la métrica escala $\approx 2{,}8$ veces conservando la mayor parte del ordenamiento ($\rho=0{,}92$) y el coeficiente de interacción se mantiene (Sección~6.5), pero $JLoss$ debe interpretarse ordinalmente y no en niveles. Tercero, once de los catorce países aportan CDS de historia larga, pero Hungría y Polonia solo tienen 14 trimestres cada uno (la ventana 2012--2015) y Pakistán apenas uno. Cuarto, Corea del Sur y Bulgaria quedan fuera del panel porque su $JLoss$ no es una medida válida a nivel de país (Sección 5.1); China, con un $JLoss$ algo elevado por un descuento de valoración menos extremo, permanece en la muestra y —por ser además el país que sostiene la significancia— es la exclusión que más condiciona el resultado. Quinto, $JLoss$ y $GaR$ son regresores estimados; la inferencia de segunda etapa que los trata como datos subestima la varianza de $\hat\theta$, si bien el error de medición atenúa el coeficiente y $\hat\theta$ resiste perturbaciones de $GaR$ de hasta el 25\,\% de su desviación. Sexto, los controles fiscales y externos exhiben signos compatibles con endogeneidad por causalidad inversa, no corregida en este diseño. Séptimo, la identificación causal del canal de \emph{nivel} (Sección~6.7) descansa en un instrumento cuya restricción de exclusión no puede verificarse directamente: el diferencial de liquidez de fondeo bancario da un efecto fuerte y significativo ($F=21{,}6$, $p=0{,}001$), pero un segundo instrumento igualmente fuerte en la primera etapa —el choque cambiario del dólar— no lo corrobora al usarse solo, y la especificación sobre-identificada con ambos \textbf{rechaza el test de Sargan} ($p=0{,}003$): los dos instrumentos no identifican el mismo parámetro, de modo que al menos uno —posiblemente ambos, en algún grado— viola la exclusión. Se reporta el de mejor argumento de exclusión (opera específicamente sobre el costo de fondeo bancario, no sobre el sovereign directamente) como el más creíble, pero el fallo de Sargan es una advertencia genuina contra tratar el efecto de nivel como causalmente cerrado. Octavo, la significancia puntual de $\theta$ es marginal ($p\approx0{,}05$ bajo las tres formas de inferencia principales), de modo que la contribución defendible es el signo y la forma del mecanismo, no la magnitud puntual.

\subsection{Agenda futura}

De las limitaciones anteriores se desprende una agenda concreta: (i) ampliar el corte transversal con economías que aporten variación de $JLoss$ —extendiendo el CDS soberano de las que hoy solo tienen $JLoss$ y $GaR$ (India vía el bono a diez años contra el UST10Y; Egipto y Pakistán a medida que su serie de CDS se completa) y cerrando el hueco de $GaR$ para Argentina, Egipto y Rusia—, que es la vía directa para que la identificación deje de descansar en tres o cuatro países; (ii) extraer de Bloomberg los componentes reales del FCI (índice accionario, REER, rendimiento 10Y) para anclar también la métrica de $GaR$ en una fuente homogénea; (iii) un \textit{bootstrap} que propague el error de estimación de la primera etapa de $GaR$ hacia la inferencia sobre $\theta$; (iv) reemplazar los datos institucionales provisionales por los \textit{Worldwide Governance Indicators} oficiales; y (v) reconciliar el rechazo del test de sobre-identificación entre el instrumento de liquidez de fondeo bancario y el del choque cambiario del dólar —ambos con primera etapa fuerte— mediante instrumentos adicionales (\textit{shocks} de política monetaria identificados, términos de intercambio por exposición sectorial) que permitan triangular cuál de los dos —o ninguno— satisface la restricción de exclusión.

\subsection{Conclusión del capítulo}

Este capítulo articula, en un único marco empírico, dos literaturas que la investigación previa había mantenido separadas: la del nexo banca-soberano y la del riesgo a la baja del crecimiento. Sobre un único panel de catorce economías emergentes construido íntegramente con datos de Bloomberg —$JLoss$ a partir de balances y precios accionarios homogéneos, spread soberano medido exclusivamente por el CDS a cinco años— la evidencia respalda el \textbf{signo y la forma} de la hipótesis central: la fragilidad bancaria sistémica y el riesgo de cola del crecimiento no determinan el spread soberano de forma aditiva e independiente, sino complementaria; el efecto de la fragilidad crece a medida que el riesgo de cola se agrava y un modelo de umbral corrobora la no linealidad sin imponer su forma. La \textbf{magnitud puntual} de la interacción es marginalmente significativa; se trata, además, de una regularidad del período posterior a la crisis financiera global cuya identificación descansa en el reducido grupo de países con variación sustantiva de la fragilidad. Estas delimitaciones precisan el alcance actual de la evidencia sin negar el mecanismo. El Capítulo~\ref{chap:paper1} deriva esta misma complementariedad desde un modelo estructural de organización industrial bancaria y pone a prueba su predicción distintiva —la amplificación por concentración—, que \emph{no} encuentra respaldo en este panel. Todos los archivos computacionales de reproducción están en el repositorio, documentados en \texttt{1\_Codigo/Panel/bbg/} y trazados a \texttt{NUMEROS\_CANONICOS\_BBG.md}; el Anexo~\ref{chap:anexoB} tabula la procedencia de cada serie.

\section*{Referencias del capítulo}
\addcontentsline{toc}{section}{Referencias del capítulo}

\begin{thebibliography}{99}
\sloppy\small
\bibitem[Acharya et al.(2014)]{AcharyaDrechslerSchnabl2014b} Acharya, V., Drechsler, I., \& Schnabl, P. (2014). A Pyrrhic Victory? Bank Bailouts and Sovereign Credit Risk. \textit{Journal of Finance}, 69(6), 2689--2739.
\bibitem[Acharya et al.(2017)]{AcharyaEtAl2017b} Acharya, V. V., Pedersen, L. H., Philippon, T., \& Richardson, M. (2017). Measuring Systemic Risk. \textit{Review of Financial Studies}, 30(1), 2--47.
\bibitem[Acharya \& Steffen(2015)]{AcharyaSteffen2015} Acharya, V. V., \& Steffen, S. (2015). The ``Greatest'' Carry Trade Ever? Understanding Eurozone Bank Risks. \textit{Journal of Financial Economics}, 115(2), 215--236.
\bibitem[Adrian \& Brunnermeier(2016)]{AdrianBrunnermeier2016b} Adrian, T., \& Brunnermeier, M. K. (2016). CoVaR. \textit{American Economic Review}, 106(7), 1705--1741.
\bibitem[Adrian \& Boyarchenko(2012)]{AdrianBoyarchenko2012} Adrian, T., \& Boyarchenko, N. (2012). Intermediary Leverage Cycles and Financial Stability. \textit{Federal Reserve Bank of New York Staff Reports}, No.\ 567.
\bibitem[Adrian et al.(2019)]{ABG2019b} Adrian, T., Boyarchenko, N., \& Giannone, D. (2019). Vulnerable Growth. \textit{American Economic Review}, 109(4), 1263--1289.
\bibitem[Allen et al.(2012)]{AllenBaliTang2012} Allen, L., Bali, T. G., \& Tang, Y. (2012). Does Systemic Risk in the Financial Sector Predict Future Economic Downturns? \textit{Review of Financial Studies}, 25(10), 3000--3036.
\bibitem[Arellano \& Bond(1991)]{ArellanoBond1991} Arellano, M., \& Bond, S. (1991). Some Tests of Specification for Panel Data: Monte Carlo Evidence and an Application to Employment Equations. \textit{Review of Economic Studies}, 58(2), 277--297.
\bibitem[Barndorff-Nielsen \& Shephard(2002)]{BarndorffNielsenShephard2002} Barndorff-Nielsen, O. E., \& Shephard, N. (2002). Econometric Analysis of Realized Volatility and Its Use in Estimating Stochastic Volatility Models. \textit{Journal of the Royal Statistical Society B}, 64(2), 253--280.
\bibitem[BIS(2006)]{BIS2006} Bank for International Settlements (BIS). (2006). \textit{International Convergence of Capital Measurement and Capital Standards (Basel II)}. Basel.
\bibitem[Billio et al.(2012)]{BillioEtAl2012} Billio, M., Getmansky, M., Lo, A. W., \& Pelizzon, L. (2012). Econometric Measures of Connectedness and Systemic Risk in the Finance and Insurance Sectors. \textit{Journal of Financial Economics}, 104(3), 535--559.
\bibitem[Bisias et al.(2012)]{BisiasEtAl2012} Bisias, D., Flood, M., Lo, A. W., \& Valavanis, S. (2012). A Survey of Systemic Risk Analytics. \textit{Annual Review of Financial Economics}, 4, 255--296.
\bibitem[Bocola(2016)]{Bocola2016} Bocola, L. (2016). The Pass-Through of Sovereign Risk. \textit{Journal of Political Economy}, 124(4), 879--926.
\bibitem[Bolton \& Jeanne(2011)]{BoltonJeanne2011} Bolton, P., \& Jeanne, O. (2011). Sovereign Default Risk and Bank Fragility in Financially Integrated Economies. \textit{IMF Economic Review}, 59(2), 162--194.
\bibitem[Brownlees \& Engle(2016)]{BrownleesEngle2016} Brownlees, C., \& Engle, R. (2016). SRISK: A Conditional Capital Shortfall Measure of Systemic Risk. \textit{Review of Financial Studies}, 30(1), 48--79.
\bibitem[Brunnermeier et al.(2016)]{BrunnermeierEtAl2016} Brunnermeier, M. K., Langfield, S., Pagano, M., Reis, R., Van Nieuwerburgh, S., \& Vayanos, D. (2016). ESBies: Safety in the Tranches. \textit{Economic Policy}, 32(90), 175--219.
\bibitem[Brunnermeier \& Sannikov(2014)]{BrunnermeierSannikov2014} Brunnermeier, M. K., \& Sannikov, Y. (2014). A Macroeconomic Model with a Financial Sector. \textit{American Economic Review}, 104(2), 379--421.
\bibitem[Calvo et al.(1996)]{CalvoLeidermanReinhart1996} Calvo, G. A., Leiderman, L., \& Reinhart, C. M. (1996). Inflows of Capital to Developing Countries in the 1990s. \textit{Journal of Economic Perspectives}, 10(2), 123--139.
\bibitem[Cameron et al.(2008)]{CameronGelbachMiller2008} Cameron, A. C., Gelbach, J. B., \& Miller, D. L. (2008). Bootstrap-Based Improvements for Inference with Clustered Errors. \textit{Review of Economics and Statistics}, 90(3), 414--427.
\bibitem[Canay(2011)]{Canay2011} Canay, I. A. (2011). A Simple Approach to Quantile Regression for Panel Data. \textit{Econometrics Journal}, 14(3), 368--386.
\bibitem[Cavallo \& Valenzuela(2010)]{CavalloValenzuela2010} Cavallo, E. A., \& Valenzuela, P. (2010). The Determinants of Corporate Risk in Emerging Markets: An Option-Adjusted Spread Analysis. \textit{International Journal of Finance \& Economics}, 15(1), 59--74.
\bibitem[Chari et al.(2024)]{Chari2024b} Chari, A., Garcés, F., Martínez, J. F., \& Valenzuela, P. (2024). Sovereign Credit Spreads, Banking Fragility, and Global Factors. \textit{Journal of Financial Stability}, 72, 101235.
\bibitem[Chernozhukov et al.(2010)]{ChernozhukovEtAl2010} Chernozhukov, V., Fernández-Val, I., \& Galichon, A. (2010). Quantile and Probability Curves Without Crossing. \textit{Econometrica}, 78(3), 1093--1125.
\bibitem[Csonto \& Ivaschenko(2013)]{CsontoIvaschenko2013} Csonto, B., \& Ivaschenko, I. (2013). Determinants of Sovereign Bond Spreads in Emerging Markets. \textit{IMF Working Paper} 13/164.
\bibitem[Dieckmann \& Plank(2012)]{DieckmannPlank2012} Dieckmann, S., \& Plank, T. (2012). Default Risk of Advanced Economies: An Empirical Analysis of CDS during the Financial Crisis. \textit{Review of Finance}, 16(4), 903--934.
\bibitem[Driscoll \& Kraay(1998)]{DriscollKraay1998} Driscoll, J. C., \& Kraay, A. C. (1998). Consistent Covariance Matrix Estimation with Spatially Dependent Panel Data. \textit{Review of Economics and Statistics}, 80(4), 549--560.
\bibitem[Edwards(1984)]{Edwards1984} Edwards, S. (1984). LDC Foreign Borrowing and Default Risk: An Empirical Investigation, 1976--80. \textit{American Economic Review}, 74(4), 726--734.
\bibitem[Farhi \& Tirole(2018)]{FarhiTirole2018b} Farhi, E., \& Tirole, J. (2018). Deadly Embrace: Sovereign and Financial Balance Sheets Doom Loops. \textit{Review of Economic Studies}, 85(3), 1781--1823.
\bibitem[Fratzscher \& Rieth(2019)]{FratzscherRieth2019} Fratzscher, M., \& Rieth, M. (2019). Monetary Policy, Bank Bailouts and the Sovereign-Bank Risk Nexus in the Euro Area. \textit{Review of Finance}, 23(4), 745--775.
\bibitem[Gennaioli et al.(2014)]{GennaioliMartinRossi2014} Gennaioli, N., Martin, A., \& Rossi, S. (2014). Sovereign Default, Domestic Banks, and Financial Institutions. \textit{Journal of Finance}, 69(2), 819--866.
\bibitem[Giglio et al.(2016)]{GiglioKellyPruitt2016} Giglio, S., Kelly, B., \& Pruitt, S. (2016). Systemic Risk and the Macroeconomy: An Empirical Evaluation. \textit{Journal of Financial Economics}, 119(3), 457--471.
\bibitem[Goldsmith-Pinkham et al.(2020)]{GoldsmithPinkhamSorkinSwift2020} Goldsmith-Pinkham, P., Sorkin, I., \& Swift, H. (2020). Bartik Instruments: What, When, Why, and How. \textit{American Economic Review}, 110(8), 2586--2624.
\bibitem[González-Rozada \& Levy-Yeyati(2008)]{GonzalezRozadaLevyYeyati2008} González-Rozada, M., \& Levy-Yeyati, E. (2008). Global Factors and Emerging Market Spreads. \textit{Economic Journal}, 118(533), 1917--1936.
\bibitem[Hansen(1999)]{Hansen1999} Hansen, B. E. (1999). Threshold Effects in Non-Dynamic Panels: Estimation, Testing, and Inference. \textit{Journal of Econometrics}, 93(2), 345--368.
\bibitem[Hansen(2000)]{Hansen2000} Hansen, B. E. (2000). Sample Splitting and Threshold Estimation. \textit{Econometrica}, 68(3), 575--603.
\bibitem[Hatzius et al.(2010)]{HatziusEtAl2010} Hatzius, J., Hooper, P., Mishkin, F. S., Schoenholtz, K. L., \& Watson, M. W. (2010). Financial Conditions Indexes: A Fresh Look after the Financial Crisis. \textit{NBER Working Paper} 16150.
\bibitem[Hausman(1978)]{Hausman1978} Hausman, J. A. (1978). Specification Tests in Econometrics. \textit{Econometrica}, 46(6), 1251--1271.
\bibitem[He \& Krishnamurthy(2013)]{HeKrishnamurthy2013} He, Z., \& Krishnamurthy, A. (2013). Intermediary Asset Pricing. \textit{American Economic Review}, 103(2), 732--770.
\bibitem[Hilscher \& Nosbusch(2010)]{HilscherNosbusch2010} Hilscher, J., \& Nosbusch, Y. (2010). Determinants of Sovereign Risk: Macroeconomic Fundamentals and the Pricing of Sovereign Debt. \textit{Review of Finance}, 14(2), 235--262.
\bibitem[Holló et al.(2012)]{HolloKremerLoDuca2012} Holló, D., Kremer, M., \& Lo Duca, M. (2012). CISS -- A Composite Indicator of Systemic Stress in the Financial System. \textit{ECB Working Paper} 1426.
\bibitem[Huang et al.(2009)]{HuangZhouZhu2009} Huang, X., Zhou, H., \& Zhu, H. (2009). A Framework for Assessing the Systemic Risk of Major Financial Institutions. \textit{Journal of Banking \& Finance}, 33(11), 2036--2049.
\bibitem[Jordà(2005)]{Jorda2005} Jordà, Ò. (2005). Estimation and Inference of Impulse Responses by Local Projections. \textit{American Economic Review}, 95(1), 161--182.
\bibitem[Kallestrup et al.(2016)]{KallestrupLandoMurgoci2016} Kallestrup, R., Lando, D., \& Murgoci, A. (2016). Financial Sector Linkages and the Dynamics of Bank and Sovereign Credit Spreads. \textit{Journal of Empirical Finance}, 38, 374--393.
\bibitem[Kealhofer(2000)]{Kealhofer2000} Kealhofer, S. (2000). The Quantification of Credit Risk. \textit{KMV Corporation} (mimeo).
\bibitem[Koenker(2004)]{Koenker2004} Koenker, R. (2004). Quantile Regression for Longitudinal Data. \textit{Journal of Multivariate Analysis}, 91(1), 74--89.
\bibitem[Koenker(2005)]{Koenker2005} Koenker, R. (2005). \textit{Quantile Regression}. Econometric Society Monographs. Cambridge University Press.
\bibitem[Koenker \& Bassett(1978)]{KoenkerBassett1978} Koenker, R., \& Bassett, G. (1978). Regression Quantiles. \textit{Econometrica}, 46(1), 33--50.
\bibitem[Laeven \& Valencia(2013)]{LaevenValencia2013} Laeven, L., \& Valencia, F. (2013). Systemic Banking Crises Database. \textit{IMF Economic Review}, 61(2), 225--270.
\bibitem[Laeven \& Valencia(2020)]{LaevenValencia2020} Laeven, L., \& Valencia, F. (2020). Systemic Banking Crises Database II. \textit{IMF Economic Review}, 68(2), 307--361.
\bibitem[Leonello(2018)]{Leonello2018} Leonello, A. (2018). Government Guarantees and the Two-Way Feedback between Banking and Sovereign Debt Crises. \textit{Journal of Financial Economics}, 130(3), 592--619.
\bibitem[Longstaff et al.(2011)]{LongstaffEtAl2011} Longstaff, F. A., Pan, J., Pedersen, L. H., \& Singleton, K. J. (2011). How Sovereign Is Sovereign Credit Risk? \textit{American Economic Journal: Macroeconomics}, 3(2), 75--103.
\bibitem[Machado \& Santos Silva(2019)]{MachadoSantosSilva2019} Machado, J. A. F., \& Santos Silva, J. M. C. (2019). Quantiles via Moments. \textit{Journal of Econometrics}, 213(1), 145--173.
\bibitem[Martin et al.(2001)]{MartinThompsonBrowne2001} Martin, R., Thompson, K., \& Browne, C. (2001). Taking to the Saddle. \textit{Risk}, 14(6), 91--94.
\bibitem[Martínez-Miera \& Repullo(2010)]{MMR2010b} Martínez-Miera, D., \& Repullo, R. (2010). Does Competition Reduce the Risk of Bank Failure? \textit{Review of Financial Studies}, 23(10), 3638--3664.
\bibitem[Merton(1974)]{Merton1974b} Merton, R. C. (1974). On the Pricing of Corporate Debt. \textit{Journal of Finance}, 29(2), 449--470.
\bibitem[Miranda-Agrippino \& Rey(2022)]{MirandaAgrippinoRey2022} Miranda-Agrippino, S., \& Rey, H. (2022). The Global Financial Cycle. \textit{Handbook of International Economics}, Vol.\ 6.
\bibitem[Mody \& Sandri(2012)]{ModySandri2012} Mody, A., \& Sandri, D. (2012). The Eurozone Crisis: How Banks and Sovereigns Came to Be Joined at the Hip. \textit{Economic Policy}, 27(70), 199--230.
\bibitem[Nickell(1981)]{Nickell1981} Nickell, S. (1981). Biases in Dynamic Models with Fixed Effects. \textit{Econometrica}, 49(6), 1417--1426.
\bibitem[Ossandon Busch et al.(2022)]{OssandonBusch2022} Ossandon Busch, M., Sánchez-Martínez, J. M., Rodríguez-Martínez, A., Montañez-Enríquez, R., \& Martínez-Jaramillo, S. (2022). Growth at Risk: Methodology and Applications in an Open-Source Platform. \textit{Latin American Journal of Central Banking}, 3, 100068.
\bibitem[Pesaran(2021)]{Pesaran2021} Pesaran, M. H. (2021). General Diagnostic Tests for Cross-Sectional Dependence in Panels. \textit{Empirical Economics}, 60(1), 13--50.
\bibitem[Prasad et al.(2019)]{PrasadEtAl2019} Prasad, A., et al. (2019). Growth at Risk: Concept and Application in IMF Country Surveillance. \textit{IMF Working Paper} 2019/036.
\bibitem[Reinhart \& Rogoff(2009)]{ReinhartRogoff2009} Reinhart, C. M., \& Rogoff, K. S. (2009). \textit{This Time Is Different: Eight Centuries of Financial Folly}. Princeton University Press.
\bibitem[Reinhart \& Rogoff(2011)]{ReinhartRogoff2011} Reinhart, C. M., \& Rogoff, K. S. (2011). From Financial Crash to Debt Crisis. \textit{American Economic Review}, 101(5), 1676--1706.
\bibitem[Remolona et al.(2008)]{RemolonaScatignaWu2008} Remolona, E., Scatigna, M., \& Wu, E. (2008). The Dynamic Pricing of Sovereign Risk in Emerging Markets: Fundamentals and Risk Aversion. \textit{Journal of Fixed Income}, 17(4), 57--71.
\bibitem[Segoviano \& Goodhart(2009)]{SegovianoGoodhart2009} Segoviano, M., \& Goodhart, C. (2009). Banking Stability Measures. \textit{IMF Working Paper} 09/4.
\bibitem[Stein(2012)]{Stein2012} Stein, J. C. (2012). Monetary Policy as Financial Stability Regulation. \textit{Quarterly Journal of Economics}, 127(1), 57--95.
\bibitem[Uribe \& Yue(2006)]{UribeYue2006} Uribe, M., \& Yue, V. Z. (2006). Country Spreads and Emerging Countries: Who Drives Whom? \textit{Journal of International Economics}, 69(1), 6--36.
\bibitem[Vasicek(1997)]{Vasicek1997} Vasicek, O. (1997). The Loan Loss Distribution. \textit{KMV Corporation}.
\bibitem[Wooldridge(2010)]{Wooldridge2010} Wooldridge, J. M. (2010). \textit{Econometric Analysis of Cross Section and Panel Data} (2nd ed.). MIT Press.
\end{thebibliography}
